pa%%%%%% NeurIPS 2026 SUBMISSION %%%%%%%%
\documentclass{article}

\usepackage{neurips_2024}            % NeurIPS style (local fallback; the
                                     % official neurips_2026.sty, if
                                     % available, will take precedence).
\usepackage[utf8]{inputenc}
\usepackage[T1]{fontenc}
\usepackage{microtype}
\usepackage{graphicx}
\usepackage{subcaption}
\usepackage{booktabs}
\usepackage{hyperref}
\usepackage{amsmath}
\usepackage{amssymb}
\usepackage{mathtools}
\usepackage{amsthm}
\usepackage[capitalize,noabbrev]{cleveref}
\usepackage{xspace}
\usepackage{multirow}
\usepackage{float}
\usepackage{algorithm}
\usepackage{algpseudocode}
\usepackage{enumitem}
\usepackage{placeins}

% Relax float placement parameters to let figures fit more easily,
% reducing end-of-body float pileup.
\renewcommand{\topfraction}{0.9}
\renewcommand{\bottomfraction}{0.8}
\renewcommand{\textfraction}{0.1}
\renewcommand{\floatpagefraction}{0.7}
\setcounter{topnumber}{3}
\setcounter{bottomnumber}{2}
\setcounter{totalnumber}{5}

%%%%%%%%%%%%%%%%%%%%%%%%%%%%%%%%
% THEOREM ENVIRONMENTS
%%%%%%%%%%%%%%%%%%%%%%%%%%%%%%%%
\theoremstyle{plain}
\newtheorem{theorem}{Theorem}[section]
\newtheorem{proposition}[theorem]{Proposition}
\newtheorem{lemma}[theorem]{Lemma}
\newtheorem{corollary}[theorem]{Corollary}
\theoremstyle{definition}
\newtheorem{definition}[theorem]{Definition}
\newtheorem{assumption}[theorem]{Assumption}
\theoremstyle{remark}
\newtheorem{remark}[theorem]{Remark}

% S-IEC name: use a simple macro that is safe inside hboxes/captions and
% also inside hyperref PDF-string contexts (section titles, etc.).
\newcommand{\SIEC}{\texorpdfstring{S\mbox{-}IEC\xspace}{S-IEC }}

\title{Decoding Efficient Diffusion: Syndrome-Guided Test-Time Error Correction}

\author{%
  Anonymous Authors \\
  Paper under double-blind review
}

\begin{document}
\maketitle

\begin{abstract}
Efficient diffusion models are typically deployed under approximation mechanisms
such as quantization, feature caching, and step reduction. These mechanisms
perturb the reverse dynamics, and the induced error is recursively propagated
across timesteps. Existing test-time correction methods mitigate this
degradation but largely treat approximation error as an unstructured local
residual. In this work, we argue that deployed diffusion should instead be
viewed as decoding over a noisy computational channel. Our starting point is a
\emph{reverse-time consistency law}: in the exact diffusion process, the
posterior mean of the clean sample forms a martingale-consistent process across
timesteps. Approximate samplers violate this law, and we interpret the
resulting martingale defect as a \emph{generalized syndrome}. We propose
Syndrome-guided Test-Time Error Correction (\SIEC), a training-free
decoder that extracts syndrome signals from redundant clean-signal estimates
at adjacent timesteps and performs minimal consensus correction before the
next reverse transition. We prove that the generalized syndrome is first-order
sensitive to structure-breaking deployment error while being insensitive to
tangent perturbations that preserve diversity; one decoding round selectively
contracts the harmful normal component of the error; and cumulative
post-correction syndrome energy controls excess denoising risk and
upper-bounds the final sampling gap to the ideal sampler under standard
discrete-time comparison conditions. These results recast efficient diffusion
correction as continuous syndrome decoding and provide a reliability-theoretic
foundation for approximate generative computation.
\end{abstract}

%=============================================================================
\section{Introduction}
\label{sec:introduction}
%=============================================================================

Denoising diffusion and score-based generative models have become a standard
paradigm for high-fidelity generation \citep{ho2020ddpm,song2021sde}, but
state-of-the-art systems are rarely deployed in their ideal form. To meet
latency, memory, and energy constraints, deployed samplers typically rely on
approximation mechanisms such as low-precision quantization, activation
caching, and aggressive step reduction. These mechanisms perturb the reverse
dynamics, and the induced error is recursively fed into subsequent denoising
steps, causing cumulative drift away from the target reverse process
\citep{zhong2026iec}. Existing test-time correction methods mitigate this
phenomenon but largely treat the approximation error as an unstructured local
perturbation to be iteratively refined \citep{zhong2026iec,yao2024tac}. This
viewpoint is natural from a generative modeling perspective, yet it overlooks
a central lesson from coding theory: reliable recovery exploits structure,
redundancy, and checkable consistency relations \citep{shannon1948,gallager1962,
berrou1993turbo}.

A similar shift in perspective is emerging for diffusion models. Denoising
Diffusion Error Correction Codes (DDECC) showed that diffusion can act as an
iterative decoder for linear codes, with parity-aware conditioning and
syndrome-guided step selection \citep{choukroun2023ddecc}. In parallel,
information-theoretic analyses have clarified that diffusion is deeply
connected to Bayesian denoising, belief propagation, and I-MMSE-type
identities \citep{guo2005immse,kong2023itdiffusion,mei2025unetbp,
reeves2025itproofs}. These lines suggest that diffusion sampling should be
viewed not only as a sequence of denoising updates, but also as a
\emph{structured decoding process}. What remains missing is a syndrome-like
object for continuous diffusion under approximate inference: unlike the
classical linear-code setting, there is no explicit parity-check matrix. If
such a syndrome is to exist, it must be extracted from the internal
consistency law of the reverse process itself.

\paragraph{Our starting point.}
Let \(m_t(x_t) = \mathbb{E}[X_0 \mid X_t = x_t]\) denote the posterior mean of
the clean sample at time \(t\). Along the exact reverse diffusion chain, we
show that \(m_t(X_t)\) is \emph{martingale-consistent}: the clean-signal
estimate at time \(t\) equals, in conditional expectation, the estimate
implied by the next reverse step. Efficient samplers violate this consistency,
and we interpret the resulting defect as a generalized syndrome
\begin{equation}
s_t
=
W_t\!\left(\hat m_t(x_t)
  - \mathbb{E}_{\tilde p_{t-1\mid t}(\cdot \mid x_t)}[\hat m_{t-1}(X_{t-1})]\right),
\label{eq:intro_syndrome}
\end{equation}
where \(\hat m_t\) is the deployed clean-signal estimator and \(W_t\) is a
preconditioning operator. In the ideal process, \(s_t = 0\); under
quantization, caching, or step reduction, one typically has \(s_t \neq 0\).
This recasts test-time correction as decoding: instead of suppressing all
local deviations, one should correct precisely those components that violate
the internal consistency law of the reverse process.

\paragraph{\SIEC.}
Based on this perspective, we propose \SIEC, a training-free
decoder that, at each reverse step, computes redundant clean-signal estimates
from adjacent timesteps, extracts the syndrome, and performs a minimal
consensus update \emph{only} when the syndrome exceeds a tolerance. This
turns test-time correction into a verifier-driven procedure: a small syndrome
certifies that the approximate trajectory remains faithful to the target
reverse process, whereas a large syndrome identifies structured deviation.
This distinction separates diversity-preserving perturbations from
manifold-violating error modes and lets us analyze correction as selective
decoding rather than generic post-processing.

\paragraph{Contributions.}
\begin{itemize}[leftmargin=1.2em,itemsep=1pt,topsep=2pt]
  \item We introduce a posterior-mean martingale view of reverse diffusion and
    define a generalized syndrome that detects violations of cross-timestep
    consistency under efficient inference.
  \item We propose \SIEC (Algorithm~\ref{alg:siec}), a training-free
    test-time correction method driven by the generalized syndrome.
  \item We prove \emph{observability} and \emph{selective stabilization}:
    syndrome-guided correction contracts the harmful manifold-violating error
    mode while preserving diversity-preserving tangent modes to first order.
  \item We establish information-theoretic bounds linking cumulative syndrome
    energy to excess denoising risk and to the sampling gap between the
    corrected sampler and the target distribution.
\end{itemize}
Compared with DDECC \citep{choukroun2023ddecc}, we do not assume an explicit
code or discrete symbols; compared with IEC and TAC \citep{zhong2026iec,
yao2024tac}, we do not iteratively reduce a generic residual but correct a
specific structural defect. \SIEC is thus closer to continuous
syndrome decoding than to heuristic refinement.

%=============================================================================
\section{Related Work}
\label{sec:related_work}
%=============================================================================

\paragraph{Efficient diffusion and test-time correction.}
A broad literature accelerates diffusion sampling through fewer denoising
steps, distilled samplers, model compression, and feature reuse
\citep{salimans2022pdistill,shang2023ptqdm,ma2024deepcache,song2023consistency}.
Several methods introduce correction mechanisms during sampling:
predictor--corrector samplers reduce discretization bias under an accurate
score model \citep{song2021sde}; TAC observes that quantization error affects
different timesteps unequally and proposes a timestep-aware correction for
low-precision diffusion \citep{yao2024tac}; IEC studies a broader family of
efficient samplers and shows that inference-time approximation errors can
accumulate exponentially across timesteps, proposing an iterative refinement
that provably reduces this growth to a linear one \citep{zhong2026iec}. These
are our closest algorithmic baselines. Rather than starting from a generic
local residual to be refined, we ask whether approximate reverse diffusion
admits an internal, syndrome-like check relation to monitor and decode.

\paragraph{Diffusion, decoding, and information theory.}
DDECC bridges diffusion and coding by using diffusion as an iterative decoder
for linear codes, with syndrome-guided step selection
\citep{choukroun2023ddecc}; that setting assumes an explicit external code,
discrete signal space, and a parity-check matrix supplied in advance. Another
line interprets U-Nets as belief propagation on generative hierarchical
models, suggesting that diffusion architectures already implement redundant
hierarchical inference \citep{mei2025unetbp}. A third line studies diffusion
through information theory: I-MMSE-type identities relate denoising objectives
to data distributions \citep{guo2005immse,kong2023itdiffusion}, and
discrete-time analyses control the divergence between an approximate chain
and an ideal one through conditional-mean estimation error
\citep{reeves2025itproofs}. These works supply the theoretical language
(posterior means, MMSE, divergence control) most relevant to our paper but do
not provide an online-computable syndrome for a deployed sampler.

\paragraph{Positioning.}
Relative to efficient-diffusion correction methods, our novelty is the
\emph{definition of the error signal}: we seek a structural defect induced by
the target process, not an undifferentiated residual
\citep{yao2024tac,zhong2026iec}. Relative to DDECC, we lack an externally
supplied code and derive the check relation from continuous reverse diffusion
itself \citep{choukroun2023ddecc}. Relative to projection-based constrained
synthesis \citep{christopher2024projected}, our verifier is endogenous: it is
induced by the reverse diffusion we seek to preserve, not an exogenous
constraint. This is the sense in which \SIEC should be understood
as \emph{continuous syndrome decoding} for approximate reverse diffusion.

%=============================================================================
\section{Preliminaries}
\label{sec:preliminaries}
%=============================================================================

\paragraph{Notation and setup.}
\(X_0 \in \mathbb{R}^d\) is a clean sample from \(P_0\); \(X_t\) denotes its
corrupted version at noise level \(t\in\{0,\dots,T\}\); \(\|\cdot\|\) is the
Euclidean norm. We use the standard discrete-time variance-preserving
formulation \citep{ho2020ddpm}: \(q(x_t \mid x_0) = \mathcal{N}(\alpha_t x_0,
\sigma_t^2 I_d)\) with \(\alpha_0=1, \sigma_0=0\), typically
\(\alpha_t^2+\sigma_t^2=1\); equivalently, \(X_t = \alpha_t X_0 +
\sigma_t\varepsilon\) with \(\varepsilon \sim \mathcal{N}(0,I_d)\), so each
step is an additive Gaussian channel. The ideal reverse kernels
\(p^\star_{t-1 \mid t}\) transport \(P_T\approx\mathcal{N}(0,I_d)\) back to
\(P_0^\star\); one ideal reverse update is
\(X_{t-1}^\star = F_t^\star(X_t^\star;\omega_t)\), where \(\omega_t\) is the
sampler's step randomness (degenerate for deterministic DDIM
\citep{song2021ddim}; Gaussian or solver-induced otherwise). The same
interface covers reverse-SDE and probability-flow ODE solvers
\citep{song2021sde}. Throughout, ideal and deployed samplers are coupled by
the same initial seed and same \(\omega_1,\dots,\omega_T\).

\paragraph{Posterior mean and deployed channel.}
The central object is the posterior mean
\(m_t(x_t) = \mathbb{E}[X_0 \mid X_t = x_t]\). Networks often predict the
noise \(\hat\varepsilon_\theta(x_t,t)\); the associated clean-signal estimate
is \(\hat x_0(x_t,t) = (x_t - \sigma_t \hat\varepsilon_\theta(x_t,t))/\alpha_t\).
For \(x_0\)- or \(v\)-prediction the model output converts to the same
quantity, so our construction is parameterization-agnostic. Under
quantization, caching, or reduced-step sampling
\citep{shang2023ptqdm,ma2024deepcache,zhong2026iec}, the deployed update is
\begin{equation}
\tilde X_{t-1} = F_t(\tilde X_t;\omega_t)
= F_t^\star(\tilde X_t;\omega_t) + \delta_t(\tilde X_t;\omega_t),
\label{eq:approx_update}
\end{equation}
where \(\delta_t\) aggregates deployment-time approximation effects (the
\emph{computational channel noise}). The pathwise error
\(e_t = \tilde X_t - X_t^\star\) evolves by dynamical amplification through
\(F_t^\star\) plus local channel noise, and \(\tilde P_0\) denotes the
deployed output law.

\paragraph{Deployment objective.}
The deployed model, scheduler, and efficiency mechanism are fixed; we do
\emph{not} allow retraining, fine-tuning, quantization-aware recalibration, or
external supervision at deployment time. A test-time correction is a
measurable map \(C_t\) acting on quantities already present in the reverse
trajectory and producing a corrected trajectory \(\bar X_{T-1},\dots,\bar X_0\)
with output law \(P_0^{\mathrm{corr}}\). We aim to make
\(D_{\mathrm{KL}}(P_0^{\mathrm{corr}} \| P_0^\star)\) small subject to a bound
on \(\mathrm{NFE}_{\mathrm{extra}}\), requiring \emph{training-free}
compatibility, \emph{selectivity} (attenuate structurally harmful drift, not
all variation), and \emph{local computability}. Conditional-mean estimation
error is known to control denoising risk and the approximate-to-target
sampling gap \citep{kong2023itdiffusion,reeves2025itproofs}; this motivates
formulating correction in terms of \(\hat x_0\). All results assume standard
local Lipschitz smoothness and bounded-moment regularity of the reverse
maps, posterior means, and channel noise; the full list of regularity
conditions is deferred to Appendix~\ref{app:reg}.

%=============================================================================
\section{Posterior-Mean Martingale and Generalized Syndrome}
\label{sec:syndrome}
%=============================================================================

The relevant error variable for approximate reverse diffusion is not the raw
denoiser discrepancy but the violation of a cross-timestep consistency law
obeyed by posterior means, in line with the information-theoretic view that
conditional-mean errors control denoising risk and the approximate-to-target
sampling gap \citep{guo2005immse,kong2023itdiffusion,reeves2025itproofs}.

\subsection{Reverse-time martingale consistency}

Define \(M_t = m_t(X_t)\) for \(t=0,\dots,T\).

\begin{proposition}[Reverse-time martingale consistency]\label{prop:martingale}
For every \(t\in\{1,\dots,T\}\) and \(P_t^\star\)-a.e.\ \(x_t\),
\(\mathbb{E}[M_{t-1} \mid X_t] = M_t\), i.e.\
\(\mathbb{E}[m_{t-1}(X_{t-1}) \mid X_t = x_t] = m_t(x_t)\).
\end{proposition}

This one-step conservation law along the ideal chain is the source of our
syndrome; a full proof via the Markov property, tower property, and
Doob-martingale construction is in Appendix~\ref{app:prop_mart}.

\subsection{Check operator and martingale defect}

Define \((T_t^\star f)(x_t) = \mathbb{E}_{p^\star_{t-1 \mid t}(\cdot \mid
x_t)}[f(X_{t-1})]\), so Proposition~\ref{prop:martingale} reads
\(m_t = T_t^\star m_{t-1}\). In deployment, \(p^\star_{t-1 \mid t}\) is
replaced by \(\tilde p_{t-1 \mid t}\) (Dirac for deterministic samplers),
inducing the deployed operator \(\tilde T_t\), and \(m_t\) is replaced by
\(\hat m_t = \hat x_0(\cdot,t)\).

\begin{definition}[Martingale defect; generalized syndrome]\label{def:defect}
The martingale defect is
\(d_t(x_t) = \hat m_t(x_t) - (\tilde T_t \hat m_{t-1})(x_t)\); for a
deterministic positive semidefinite \(W_t\in\mathbb{R}^{d\times d}\), the
\emph{generalized syndrome} is
\begin{equation}
s_t(x_t) = W_t\, d_t(x_t)
 = W_t\!\left(\hat m_t(x_t) - (\tilde T_t \hat m_{t-1})(x_t)\right).
\label{eq:syndrome}
\end{equation}
\end{definition}

\begin{proposition}[Zero-syndrome property]\label{prop:zero_syn}
If \(\tilde p_{t-1 \mid t} = p^\star_{t-1 \mid t}\) and \(\hat m_t = m_t\)
for all \(t\), then \(d_t \equiv s_t \equiv 0\) \(P_t^\star\)-a.e.
\end{proposition}

The syndrome is thus zero on the ideal process, endogenous, and
MMSE-aligned; it plays the role of a parity test against the
martingale-consistent family \(\{m_t\}_t\), induced endogenously by the
reverse process rather than supplied by an external code
\citep{choukroun2023ddecc}. The \emph{cumulative syndrome energy} is
\(\mathcal{E}_{\mathrm{syn}} = \sum_{t=1}^T \mathbb{E}\|s_t(\tilde X_t)\|^2\).

\subsection{A computable one-sample proxy}

\begin{proposition}[Unbiased proxy]\label{prop:proxy}
Let \(\hat s_t^{(1)} = W_t(\hat x_0(\tilde X_t,t) - \hat x_0(\tilde
X_{t-1}^{\mathrm{tent}}, t-1))\), where \(\tilde X_{t-1}^{\mathrm{tent}}
\sim \tilde p_{t-1\mid t}(\cdot \mid \tilde X_t)\). Then
\(\mathbb{E}[\hat s_t^{(1)} \mid \tilde X_t] = s_t(\tilde X_t)\), with
equality for deterministic samplers.
\end{proposition}

The \(K\)-sample average \(\hat s_t^{(K)}\) reduces variance; \(K=1\) is
the default. Full proofs of Propositions~\ref{prop:zero_syn}--\ref{prop:proxy}
and a variance bound for \(\hat s_t^{(K)}\) are in
Appendix~\ref{app:prop_syn}.

%=============================================================================
\section{S-IEC: Syndrome-Guided Test-Time Error Correction}
\label{sec:method}
%=============================================================================

\SIEC is a training-free decoder on top of a deployed efficient sampler. At
each reverse step it performs a tentative transition, checks posterior-mean
consistency, and only if the check fails applies a minimal correction.

\paragraph{Correction interface.}
Any deployed sampler factors as \(F_t(x_t;\omega_t) = \Psi_t(x_t, \hat
x_0(x_t,t);\omega_t)\) where \(\Psi_t\) is the scheduler map (DDPM, DDIM,
reverse-SDE \citep{ho2020ddpm,song2021ddim,song2021sde}); e.g.\
\(\Psi_t^{\mathrm{DDIM}}(x,z) = \alpha_{t-1} z +
(\sigma_{t-1}/\sigma_t)(x-\alpha_t z)\). \SIEC modifies neither the model
nor the scheduler; it only replaces the clean-signal estimate supplied to
\(\Psi_t\) when the syndrome signals inconsistency.

\paragraph{Event-triggered consensus correction.}
Given \(\tilde x_t\), compute \(\hat x_{0,t}\) and tentative next state
\(\tilde x_{t-1}^{(0)} = \Psi_t(\tilde x_t,\hat x_{0,t};\omega_t)\); the
lookahead estimate \(\hat x_{0,t-1}^{(0)} = \hat x_0(\tilde
x_{t-1}^{(0)},t-1)\) yields the proxy \(\hat s_t^{(0)} = W_t(\hat x_{0,t} -
\hat x_{0,t-1}^{(0)})\) and score \(r_t^{(0)} = \tfrac1d \|\hat s_t^{(0)}\|^2\).
With threshold \(\tau_t\), the rule is: \emph{accept} \(\tilde x_{t-1}^{(0)}\)
iff \(r_t^{(0)} \le \tau_t\). Accepted steps cache \(\hat x_{0,t-1}^{(0)}\)
for the next step so the lookahead is not wasted. Otherwise a local correction
activates.

\paragraph{Consensus update in clean-estimate space.}
Because the syndrome is a defect of posterior-mean consistency, correction
happens in clean-estimate space. For weights \(\Sigma_t,\Sigma_{t-1}\succ 0\)
and gain \(\lambda_t>0\), the consensus closed form is \(\bar x_{0,t} =
(\Sigma_t^{-1} + \lambda_t\Sigma_{t-1}^{-1})^{-1}(\Sigma_t^{-1}\hat x_{0,t} +
\lambda_t\Sigma_{t-1}^{-1}\hat x_{0,t-1})\). The isotropic
\(\Sigma_t=\Sigma_{t-1}=I_d\) yields
\begin{equation}
\bar x_{0,t} = \hat x_{0,t} - \gamma_t\, \hat s_t,
\qquad
\gamma_t = \tfrac{\lambda_t}{1+\lambda_t} \in (0,1),
\label{eq:iso_consensus}
\end{equation}
so the corrected scheduler transition reuses the same \(\omega_t\):
\(\tilde x_{t-1}^{(1)} = \Psi_t(\tilde x_t, \bar x_{0,t}; \omega_t)\).
Freezing \(\omega_t\) isolates structural error from sampling noise.

\begin{algorithm}[t]
\caption{\SIEC reverse step at time \(t\)}
\label{alg:siec}
\begin{algorithmic}[1]
\Require \(\tilde x_t\); scheduler \(\Psi_t\), randomness \(\omega_t\);
  \(W_t, \tau_t, \lambda_t\); max rounds \(L_t\).
\State \(\hat x_{0,t} \gets \hat x_0(\tilde x_t,t)\);
       \(\tilde x_{t-1}^{(0)} \gets \Psi_t(\tilde x_t,\hat x_{0,t};\omega_t)\)
\For{\(\ell = 0,\dots,L_t-1\)}
  \State \(\hat x_{0,t-1}^{(\ell)} \gets \hat x_0(\tilde x_{t-1}^{(\ell)},t-1)\);
         \(\hat s_t^{(\ell)} \gets W_t(\hat x_{0,t} - \hat x_{0,t-1}^{(\ell)})\);
         \(r_t^{(\ell)} \gets \tfrac1d\|\hat s_t^{(\ell)}\|^2\)
  \If{\(r_t^{(\ell)} \le \tau_t\)}
     \Return \(\bar x_{t-1} \gets \tilde x_{t-1}^{(\ell)}\);
              cache \(\hat x_{0,t-1}^{(\ell)}\) \Comment{accept}
  \EndIf
  \State \(\bar x_{0,t}^{(\ell+1)} \gets \hat x_{0,t} - \gamma_t \hat s_t^{(\ell)}\);
         \(\tilde x_{t-1}^{(\ell+1)} \gets \Psi_t(\tilde x_t, \bar x_{0,t}^{(\ell+1)};\omega_t)\)
\EndFor
\State \Return \(\bar x_{t-1} \gets \tilde x_{t-1}^{(L_t)}\)
\end{algorithmic}
\end{algorithm}

\paragraph{Complexity and positioning.}
Algorithm~\ref{alg:siec} requires no parameter updates, calibration data,
or auxiliary network; extra compute is bounded by
\(\sum_t L_t\mathbb{P}(r_t^{(0)}>\tau_t)\) (Corollary~\ref{cor:tradeoff}).
In practice \(L_t=1\) suffices; multi-round decoding is needed only under
severe approximation (e.g., low-bit quantization, aggressive caching)
\citep{shang2023ptqdm,ma2024deepcache}. TAC and IEC correct under
timestep-dependent accumulation but lack a process-induced zero-check
\citep{yao2024tac,zhong2026iec}; DDECC decodes a code-supplied parity
\citep{choukroun2023ddecc}. \SIEC is training-free and lightweight, but
driven by a verifier derived endogenously from cross-timestep posterior-mean
consistency.

%=============================================================================
\section{Main Theoretical Results}
\label{sec:theory}
%=============================================================================

We identify the deployment-error component that is \emph{observable} through
the reverse process, show that one \SIEC round \emph{selectively contracts}
it, and connect cumulative syndrome energy to both excess denoising risk and
the final sampling gap. Results are stated for deterministic schedulers and
exact syndromes; by Proposition~\ref{prop:proxy} the expectation-level
bounds carry over to \(\hat s_t^{(1)}\). Full proofs are in
Appendix~\ref{app:thm_proofs}.

\subsection{Geometric decomposition}

Let \(x_t^\star\) be the ideal trajectory and \(\bar x_t\) the trajectory
returned by \SIEC under the same seed; define \(e_t = \bar x_t - x_t^\star\).
At the ideal trajectory let \(J_t = \nabla m_t(x_t^\star)\),
\(B_t = \nabla_x\Psi_t^\star(x_t^\star,m_t(x_t^\star))\),
\(C_t = \nabla_z\Psi_t^\star(x_t^\star,m_t(x_t^\star))\), and
\(G_t = \nabla d_t(x_t^\star)\). The local tangent and normal subspaces of
the posterior-mean level set are \(\mathcal{T}_t = \ker J_t\) and
\(\mathcal{N}_t = \mathcal{T}_t^\perp\), with orthogonal projectors
\(P_t^\parallel, P_t^\perp\). Every local error decomposes as
\(e_t = e_t^\parallel + e_t^\perp\); \(e_t^\parallel\) leaves the posterior
mean unchanged to first order, while \(e_t^\perp\) changes it to first order.

\begin{assumption}[Local regularity]\label{asm:reg}
There exists \(r_t>0\) such that on \(B(x_t^\star,r_t)\), the maps
\(m_t, \hat m_t, \Psi_t^\star, \Psi_t\) are twice continuously differentiable;
and there exist mismatch levels \(\zeta_t,\xi_t\ge 0\) with
\(\|\hat m_t - m_t\|_{C^1(B)} \le \zeta_t\) and
\(\|\Psi_t - \Psi_t^\star\|_{C^1(B)} \le \xi_t\).
\end{assumption}

\begin{assumption}[Syndrome observability]\label{asm:obs}
There exist \(0 \le \beta_t < \kappa_t\) with
\(\|W_t G_t u\| \le \beta_t\|u\|\) for \(u\in\mathcal{T}_t\) and
\(\|W_t G_t v\| \ge \kappa_t\|v\|\) for \(v\in\mathcal{N}_t\).
\end{assumption}

\subsection{Syndrome observability}

\begin{theorem}[First-order syndrome expansion]\label{thm:observability}
Under Assumptions~\ref{asm:reg}--\ref{asm:obs}, there exist \(c_{t,1},c_{t,2}>0\)
such that for every \(\bar x_t \in B(x_t^\star,r_t)\),
\begin{equation}
s_t(\bar x_t) = W_t G_t e_t + b_t + r_t,
\label{eq:syn_expansion}
\end{equation}
with \(\|b_t\|\le c_{t,1}(\xi_t+\zeta_t)\) and
\(\|r_t\|\le c_{t,2}(\|e_t\|^2 + (\xi_t+\zeta_t)\|e_t\| + (\xi_t+\zeta_t)^2)\).
Consequently,
\begin{equation}
\|e_t^\perp\|
\;\le\;
\tfrac{1}{\kappa_t}\left(\|s_t(\bar x_t)\| + \beta_t\|e_t^\parallel\|
  + \|b_t\| + \|r_t\|\right).
\label{eq:obs_bound}
\end{equation}
\end{theorem}

\emph{Proof sketch.} Taylor-expand \(d_t\) at \(x_t^\star\); the leading
term is \(G_t e_t\), estimator/scheduler mismatch contribute \(b_t\), and
the quadratic remainder is \(r_t\). Projecting onto
\(\mathcal{T}_t\oplus\mathcal{N}_t\) and applying Assumption~\ref{asm:obs}
gives~\eqref{eq:obs_bound}; see Appendix~\ref{app:thm_obs}. When
\(\beta_t\) is small, a small syndrome certifies small \(e_t^\perp\) even if
\(e_t^\parallel\) remains nonzero---the syndrome is a selective verifier of
structurally damaging deviation.

\subsection{Selective stabilization under one decoding round}

In the isotropic regime \(W_t=\Sigma_t=\Sigma_{t-1}=I_d\), the consensus
gain is \(\gamma_t=\lambda_t/(1+\lambda_t)\), the corrected clean estimate
is \(\bar x_{0,t} = \hat x_{0,t} - \gamma_t\hat s_t\), and the linearized
one-step corrected propagation is governed by
\begin{equation}
L_t(\gamma_t) = B_t + C_t J_t - \gamma_t\, C_t G_t.
\label{eq:L_operator}
\end{equation}

\begin{assumption}[Selective decodability]\label{asm:decode}
There exist \(\rho_t\in(0,1)\), \(\nu_t,\chi_t\ge 0\) with
\(\|P_{t-1}^\perp L_t(\gamma_t) P_t^\perp\|\le \rho_t\),
\(\|P_{t-1}^\parallel L_t(\gamma_t) P_t^\parallel\|\le 1+\nu_t\), and
\(\|P_{t-1}^\perp L_t(\gamma_t) P_t^\parallel\|
 +\|P_{t-1}^\parallel L_t(\gamma_t) P_t^\perp\|\le \chi_t\).
\end{assumption}

\begin{theorem}[Selective stabilization]\label{thm:stabilization}
Under Assumptions~\ref{asm:reg}--\ref{asm:decode}, there exist
\(c_{t,3},c_{t,4}>0\) such that one \SIEC round yields
\begin{align}
\|P_{t-1}^\perp e_{t-1}^{\mathrm{corr}}\| &\le \rho_t\|e_t^\perp\|
  + \chi_t\|e_t^\parallel\| + c_{t,3}(\xi_t+\zeta_t+\|e_t\|^2),
  \label{eq:stab_normal}\\
\|P_{t-1}^\parallel e_{t-1}^{\mathrm{corr}}\| &\le (1+\nu_t)\|e_t^\parallel\|
  + \chi_t\|e_t^\perp\| + c_{t,4}(\xi_t+\zeta_t+\|e_t\|^2).
  \label{eq:stab_tangent}
\end{align}
\end{theorem}

\emph{Proof sketch.} Linearize \(\bar x_{t-1} = \Psi_t(\bar x_t,\bar
x_{0,t};\omega_t)\) around the ideal trajectory: the consensus update
injects \(-\gamma_t C_t G_t e_t\)---exactly the term missing from the
uncorrected dynamics---and Assumption~\ref{asm:decode} gives contraction on
\(\mathcal{N}_t\) with weak crosstalk \(\chi_t\); see
Appendix~\ref{app:thm_stab}. Without the correction term the normal
contraction fails, recovering the accumulation regime of
\citep{zhong2026iec}.

\subsection{From syndrome energy to excess MMSE and KL gap}

Let \(\bar m_t, \bar T_t\) denote the clean-signal estimator and one-step
transition operator of the corrected sampler, and define the
post-correction syndrome \(\bar s_t = W_t(\bar m_t - \bar T_t \bar m_{t-1})\)
and the stepwise excess MMSE
\(\Delta_t = \mathbb{E}\|\bar m_t(\bar X_t) - m_t(\bar X_t)\|^2\),
\(\Delta_0 = 0\). The orthogonality principle gives that \(\Delta_t\) is
exactly the excess denoising risk at step \(t\)
\citep{guo2005immse,kong2023itdiffusion}.

\begin{assumption}[Operator stability]\label{asm:op}
There exists \(a\in[0,1)\) with
\(\mathbb{E}\|(\bar T_t f)(\bar X_t)\|^2 \le a\,\mathbb{E}\|f(\bar X_{t-1})\|^2\)
for all measurable \(f\); and constants \(0<\underline w\le\overline w<\infty\)
with \(\underline w \le \lambda_{\min}(W_t^\top W_t) \le
\lambda_{\max}(W_t^\top W_t) \le \overline w\).
\end{assumption}

\begin{theorem}[Excess MMSE from syndrome energy]\label{thm:mmse}
Under Assumption~\ref{asm:op}, let
\(\eta_t^2 = \mathbb{E}\|((\bar T_t - T_t^\star) m_{t-1})(\bar X_t)\|^2\). Then
\begin{equation}
\Delta_t \;\le\; \tfrac{3}{\underline w}\,\mathbb{E}\|\bar s_t(\bar X_t)\|^2
 + 3 a \Delta_{t-1} + 3\eta_t^2,
\label{eq:mmse_recursion}
\end{equation}
and with
\(\mathcal{E}^{\mathrm{corr}}_{\mathrm{syn}} = \sum_{t=1}^T \mathbb{E}\|\bar
s_t(\bar X_t)\|^2\),
\begin{equation}
\sum_{t=1}^T \Delta_t
\;\le\;
\tfrac{3}{(1-a)\underline w}\,\mathcal{E}^{\mathrm{corr}}_{\mathrm{syn}}
 + \tfrac{3}{1-a}\sum_{t=1}^T \eta_t^2.
\label{eq:mmse_sum}
\end{equation}
\end{theorem}

\emph{Proof sketch.} Apply \(\bar m_t - m_t = \bar d_t + \bar T_t(\bar
m_{t-1} - m_{t-1}) + (\bar T_t - T_t^\star)m_{t-1}\); squared norms with
\(\|u+v+w\|^2 \le 3(\|u\|^2+\|v\|^2+\|w\|^2)\), Assumption~\ref{asm:op}, and
weight conditioning; unroll. See Appendix~\ref{app:thm_mmse}.

\begin{theorem}[Distributional control]\label{thm:kl}
Under the schedule regularity of \citep{reeves2025itproofs}
(Appendix~\ref{app:reeves}), there exists \(C_{\mathrm{diff}}>0\) depending
only on the schedule with
\begin{equation}
D_{\mathrm{KL}}(P_0^{\mathrm{corr}} \,\|\, P_0^\star)
\;\le\;
C_{\mathrm{diff}}\!\sum_{t=1}^T \Delta_t
\;\le\;
\tfrac{3C_{\mathrm{diff}}}{(1-a)\underline w}\,\mathcal{E}^{\mathrm{corr}}_{\mathrm{syn}}
 + \tfrac{3C_{\mathrm{diff}}}{1-a}\sum_{t=1}^T \eta_t^2.
\label{eq:kl_bound}
\end{equation}
By Pinsker, \(\|P_0^{\mathrm{corr}} - P_0^\star\|_{\mathrm{TV}} \le
\sqrt{\tfrac12 D_{\mathrm{KL}}}\): cumulative syndrome energy controls
total-variation error.
\end{theorem}

\emph{Proof sketch.} The first inequality is the discrete-time comparison
of \citep{reeves2025itproofs} specialized to the corrected chain; the
second follows from Theorem~\ref{thm:mmse}. See Appendix~\ref{app:thm_kl}.
Together, Theorems~\ref{thm:mmse}--\ref{thm:kl} make the syndrome a
\emph{global accounting variable}.

\subsection{Event-triggered trade-off}

\begin{corollary}[Event-triggered trade-off]\label{cor:tradeoff}
If whenever correction activates,
\(\mathbb{E}[\|\bar s_t\|^2 \mid r_t^{(0)} > \tau_t] \le q_t\,\mathbb{E}[\|\hat
s_t^{(0)}\|^2 \mid r_t^{(0)} > \tau_t]\) with \(q_t\in(0,1)\), then
\begin{equation}
\mathcal{E}^{\mathrm{corr}}_{\mathrm{syn}} \;\le\;
d\sum_{t=1}^T \tau_t
+ \sum_{t=1}^T q_t\,\mathbb{E}\bigl[\|\hat s_t^{(0)}\|^2
  \mathbf{1}\{r_t^{(0)}>\tau_t\}\bigr],
\label{eq:triggered_energy}
\end{equation}
and \(\mathbb{E}[\mathrm{NFE}_{\mathrm{extra}}] \le
\sum_t L_t\,\mathbb{P}(r_t^{(0)} > \tau_t)\).
\end{corollary}

Substituting into~\eqref{eq:kl_bound} gives an explicit upper bound on the
final KL gap as a function of trigger thresholds, activation frequency, and
per-step syndrome reduction; see Appendix~\ref{app:cor_tradeoff}.

\section{Experiments: Theory Validation}
\label{sec:experiments}
%=============================================================================

We validate four predictions under a paired coupling protocol (same initial
seed, same \(\omega_t\) for ideal and deployed samplers): \emph{(i)} the
syndrome is an observable of the harmful component
(Theorem~\ref{thm:observability}); \emph{(ii)} one \SIEC round selectively
contracts the unstable component (Theorem~\ref{thm:stabilization});
\emph{(iii)} cumulative syndrome energy predicts excess MMSE and the
sampling gap (Theorems~\ref{thm:mmse}--\ref{thm:kl}); \emph{(iv)}
event-triggered correction yields a compute--quality trade-off
(Corollary~\ref{cor:tradeoff}).

\paragraph{Setup.}
We use a low-rank Gaussian model \(X_0 = U Z\), \(Z \sim \mathcal{N}(0,\Lambda)\),
\(U \in \mathbb{R}^{d\times k}\), \(U^\top U = I_k\), with \(d=1024\),
\(k=32\), \(T=200\), cosine schedule, and \(\lambda_i \in [0.1,10]\) at
log-spacing (condition number 100). The posterior mean is explicit
(Appendix~\ref{app:exp_details}) and
\(\mathcal{T}_t = \mathrm{span}(U)^\perp\), \(\mathcal{N}_t = \mathrm{span}(U)\),
giving exact access to \(e_t^{\parallel/\perp}, \Delta_t, D_{\mathrm{KL}}\).
Deployment error is injected via \(\hat m_t = m_t + \eta\, n\),
\(n\sim\mathcal{N}(0,I_d)\), with variants for tangent/normal-projected
noise, simulated 8-bit/4-bit uniform quantization, and SNR-dependent
\(\eta_t = \eta_{\mathrm{base}}\sigma_t^2/(\alpha_t^2+\sigma_t^2)\). \SIEC
uses \(W_t=I_d\), \(\lambda_t = c\sigma_t^2/(\alpha_t^2+\sigma_t^2)\) with
\(c=1\), \(L_t=1\), and thresholds \(\tau_t\) calibrated from 100--128 pilot
trajectories. We compare against uncorrected, always-on \SIEC,
event-triggered \SIEC at percentiles \(\{50,70,80,90,95\}\), random-trigger
(matched budget), uniform-periodic, and naive always-on refinement.

\paragraph{Observability (Theorem~\ref{thm:observability}).}
Equal-norm perturbations along \(\mathcal{T}_t\) or \(\mathcal{N}_t\) at five
timesteps (80 trajectories each) confirm: \(\|s_t\|\) is linear in
\(\|e_t^\perp\|\) (slope \(\approx\kappa_t\)) and essentially flat in
\(\|e_t^\parallel\|\); the empirical ratio \(\kappa_t/\beta_t >50\) at all
tests, strongly satisfying Assumption~\ref{asm:obs}
(Figure~\ref{fig:observability}).

\begin{figure}[t]
\centering
\includegraphics[width=0.65\linewidth]{fig1_observability.png}
\caption{\textbf{Observability.} (a) \(\|s_t\|\) vs.\ \(\|e_t^\perp\|\):
linear, slope \(\approx\kappa_t\). (b) vs.\ \(\|e_t^\parallel\|\):
near-flat. (c) \(\kappa_t/\beta_t >50\).}
\label{fig:observability}
\end{figure}

\paragraph{Selective stabilization (Theorem~\ref{thm:stabilization}).}
Sweeping \(\eta\in\{0.02,0.05,0.1,0.2\}\), \(\lambda\in\{0.1,\dots,5\}\)
(100 trajectories per setting), \SIEC reduces error in both components,
with mean \(\hat\rho^\perp = 0.82\), \(\hat\rho^\parallel=0.79\); normal
contraction improves monotonically in \(\lambda\) and saturates near
\(\lambda\approx 2\) (Figure~\ref{fig:stabilization}).

\begin{figure}[t]
\centering
\includegraphics[width=0.65\linewidth]{fig2_stabilization.png}
\caption{\textbf{Stabilization.} Contraction ratios
\(\hat\rho^\perp, \hat\rho^\parallel\) across all \((\eta,\lambda,t)\);
both below~1. Right panels: dependence on \(\lambda\).}
\label{fig:stabilization}
\end{figure}

\paragraph{Energy controls MMSE and KL
(Theorems~\ref{thm:mmse}--\ref{thm:kl}).}
Sweeping \(\eta \in [10^{-2},10^{-0.3}]\), \(c\in\{0,0.5,1,2\}\), and
\(\tau\in\{70,90,100\}\) yields 60+ configurations (100 trajectories each):
cumulative syndrome energy is monotonically related to both
\(\sum_t\Delta_t\) and \(D_{\mathrm{KL}}\), validating
\eqref{eq:mmse_sum}--\eqref{eq:kl_bound} (Figure~\ref{fig:mmse_kl} in
Appendix~\ref{app:exp_details}).

\paragraph{Compute--quality trade-off (Corollary~\ref{cor:tradeoff}).}
At \(\eta=0.3\), event-triggered \SIEC traces a smooth Pareto front in
\(\tau\) (Figure~\ref{fig:pareto} and Table~\ref{tab:pareto} in
Appendix~\ref{app:exp_details}): at matched NFE \(\sim 240\),
\SIEC(\(\tau\)=80th) attains final MSE \(0.0913\), beating random-trigger
(\(0.0919\)) and uniform-periodic (\(0.0918\)); always-on correction cuts
path error by \(\sim\!12\%\). In learned diffusion where errors concentrate
at specific timesteps, the selective advantage is expected to be more
pronounced.

\paragraph{Robustness.}
Across six deployment scenarios (mild/moderate/severe isotropic, 8-/4-bit
quantization, SNR-dependent), \SIEC consistently reduces per-step error in
proportion to perturbation severity and stabilizes the full trajectory,
not just the final step (Figure~\ref{fig:trajectories} in
Appendix~\ref{app:exp_details}).

\section{Discussion and Limitations}
\label{sec:discussion}
%=============================================================================

\SIEC reinterprets test-time correction as a decoding problem. Its control
signal is a verifier derived endogenously from cross-timestep posterior-mean
consistency, placing it between DDECC (which decodes an externally supplied
code \citep{choukroun2023ddecc}) and IEC-style refinement (which reduces an
undifferentiated residual \citep{zhong2026iec}); the redundancy comes from
adjacent clean-signal beliefs, and the verifier is built from the same
conditional-mean objects that govern denoising risk
\citep{mei2025unetbp,kong2023itdiffusion,reeves2025itproofs}.

\paragraph{Limitations.}
\begin{itemize}[leftmargin=1.2em,itemsep=1pt,topsep=2pt]
  \item \emph{Local guarantees.} Theorems~\ref{thm:observability}--\ref{thm:kl}
    rely on local smoothness and observability near the ideal trajectory.
  \item \emph{Estimator quality.} A severely mismatched \(\hat x_0\) mixes
    deployment drift with denoiser bias; our mismatch terms separate
    these only perturbatively.
  \item \emph{Multistep redundancy.} Adjacent-step agreement weakens under
    heavy distillation \citep{song2023consistency}; multi-branch or
    multi-scale consistency is a natural extension.
  \item \emph{Internal vs.\ semantic.} A small syndrome certifies internal
    self-consistency, not prompt fidelity or safety; external verifiers
    \citep{christopher2024projected} are complementary.
  \item \emph{Average-case, not adversarial.} Our bounds control expected
    bounded error, not worst-case adversarial perturbations.
\end{itemize}

\paragraph{Outlook.}
Three directions stand out: \emph{(i)} an explicit \emph{computational channel
model} with admissible fault classes and bounded-distance decoding, making
the decodable approximation radius explicit; \emph{(ii)} richer syndromes
across scales and skip states \citep{mei2025unetbp}; \emph{(iii)}
\emph{adaptive protection}---unequal-error-protection for trajectories,
timesteps, or semantic regions flagged by the syndrome.

%=============================================================================
\section{Conclusion}
\label{sec:conclusion}
%=============================================================================

We identified posterior-mean martingale consistency as the internal
conservation law of the exact reverse process and defined its violation as
a generalized syndrome. Based on this object we proposed \SIEC, a
training-free decoder that performs syndrome-guided consensus correction
using redundant clean-signal estimates already present along the reverse
trajectory: the syndrome is a first-order observable of the harmful
component of deployment drift, one round selectively contracts it, and
cumulative syndrome energy controls both excess denoising risk and the
final sampling gap. The missing object in efficient diffusion was not
another denoiser but an \emph{endogenous syndrome} induced by the target
process itself.

\section*{Impact Statement}

\SIEC enables test-time error correction without retraining, lowering the
computational cost, energy use, and carbon emissions of large-scale
diffusion inference.

\bibliographystyle{plainnat}
\bibliography{mybib}

\clearpage
\appendix

%=============================================================================
%                              APPENDIX
%=============================================================================

\section{Regularity Conditions and Comparison Lemma}
\label{app:reg}

We collect here the regularity conditions referred to in the main text and
restate the discrete-time comparison lemma of
\citet{reeves2025itproofs} in a form suitable for our use.

\subsection{Standing regularity conditions}
\label{app:regconds}

Throughout the paper we assume that the deployed model parameters and
scheduler are fixed, and we compare the ideal and deployed samplers under
the \emph{common-randomness coupling}: \(X_T^\star = \tilde X_T = \bar X_T
\sim \mathcal{N}(0,I_d)\) and the same \(\omega_1,\dots,\omega_T\). Under
this coupling, we impose the following four conditions.

\paragraph{(R1) Local smoothness of reverse dynamics.}
For each \(t\in\{1,\dots,T\}\), there exists a radius \(r_t>0\) such that on
the Euclidean ball \(B(x_t^\star, r_t)\) centered at the ideal trajectory,
the ideal reverse map \(F_t^\star\), the ideal scheduler \(\Psi_t^\star\),
the deployed scheduler \(\Psi_t\), and the posterior mean \(m_t\) are
twice continuously differentiable with uniformly bounded second derivatives.

\paragraph{(R2) Bounded-moment channel noise.}
The computational channel noise \(\delta_t\) in~\eqref{eq:approx_update}
satisfies
\begin{equation}
\mathbb{E}\!\left[\|\delta_t(X_t^\star;\omega_t)\|^2 \mid X_t^\star\right]
\;\le\; \eta_t^2
\qquad\text{for all }t\in\{1,\dots,T\},
\label{eq:app_bounded_noise}
\end{equation}
with \(\eta_t\ge 0\) deterministic.

\paragraph{(R3) Coupled trajectory remains local.}
With probability one, the coupled approximate trajectory
\(\{\bar X_t\}_{t=T}^{0}\) stays in the local balls,
\(\bar X_t \in B(x_t^\star, r_t)\) for all \(t\). This is the regime in
which first-order linearization of the reverse dynamics is informative.

\paragraph{(R4) Scheduler regularity.}
\(\alpha_t>0\) for all \(t<T\), and the noise levels \(\sigma_t\) are
monotone in \(t\). In addition, \(\alpha_T=0\) or \(\sigma_T\ge 1-\epsilon\)
for some small \(\epsilon>0\), so that the starting distribution
\(P_T\) is close to \(\mathcal{N}(0,I_d)\).

Conditions (R1)--(R4) are standard for discrete-time stability analyses of
diffusion sampling \citep{reeves2025itproofs,zhong2026iec} and are
satisfied, in particular, by cosine and linear schedules on the
low-rank Gaussian model used in our experiments.

\subsection{Mismatch levels.}
Assumption~\ref{asm:reg} in the main text introduces nonnegative mismatch
levels \(\zeta_t\) and \(\xi_t\); for completeness, these are
\begin{align*}
\zeta_t &\;\ge\;
  \sup_{x\in B(x_t^\star,r_t)}
  \left[\|\hat m_t(x) - m_t(x)\| + \|\nabla \hat m_t(x) - \nabla m_t(x)\|\right],\\
\xi_t &\;\ge\;
  \sup_{(x,z)\in \mathcal{U}_t}
  \left[\|\Psi_t(x,z;\omega_t) - \Psi_t^\star(x,z;\omega_t)\|
   + \|\nabla \Psi_t(x,z;\omega_t) - \nabla \Psi_t^\star(x,z;\omega_t)\|\right],
\end{align*}
where \(\mathcal{U}_t\) is the local neighborhood of
\((x_t^\star, m_t(x_t^\star))\) in the product space. These bounds are
deterministic but allowed to be time-varying.

\subsection{Discrete-time comparison lemma}
\label{app:reeves}

We restate \citet[Theorem~2]{reeves2025itproofs} in the specific form used
to prove Theorem~\ref{thm:kl}. Let \(\bar P\) and \(P^\star\) be the laws on
\((\mathbb{R}^d)^{T+1}\) of the corrected and ideal reverse chains
respectively, sharing the initial distribution \(\mathcal{N}(0,I_d)\) at
\(t=T\) and the same noise schedule.

\begin{lemma}[Discrete-time comparison]\label{lem:reeves}
Under (R1)--(R4), there exists a constant \(C_{\mathrm{diff}}>0\) depending
only on the noise schedule \(\{(\alpha_t,\sigma_t)\}_{t=1}^T\) such that the
marginals \(\bar P_0\) and \(P_0^\star\) at \(t=0\) satisfy
\begin{equation}
D_{\mathrm{KL}}\!\left(\bar P_0 \,\|\, P_0^\star\right)
\;\le\;
C_{\mathrm{diff}} \sum_{t=1}^T
\mathbb{E}\!\left[\,\|\bar m_t(\bar X_t) - m_t(\bar X_t)\|^2\,\right],
\label{eq:lem_reeves}
\end{equation}
where \(\bar m_t\) is the clean-signal estimator used by the corrected
chain at step \(t\), \(m_t\) is the ideal posterior mean, and the
expectation is under \(\bar P\).
\end{lemma}

\begin{proof}[Proof sketch]
The full derivation is in \citet{reeves2025itproofs}; we only recall the
mechanism. First, the KL divergence between the two reverse chains decomposes
as a sum of per-step conditional KL terms using the chain rule of relative
entropy. Second, under (R1) and (R4) each per-step conditional KL is
controlled by the squared \(L^2\) distance between the means of the two
Gaussian (or Gaussian-limit) one-step transition kernels, which itself is
controlled by the squared error in the clean-signal estimate. Third, the
scheduler coefficients are absorbed into the single constant
\(C_{\mathrm{diff}}\), which for the variance-preserving discretization
takes the explicit form
\(C_{\mathrm{diff}} = \sup_{t}\alpha_t^2/(2\sigma_t^2)\) up to a
multiplicative constant depending on the schedule regularity window; see
\citet[eq.~(14)]{reeves2025itproofs}.
\end{proof}

Under (R4), \(C_{\mathrm{diff}}\) is finite for standard cosine and linear
schedules. In the main text we write
\(\Delta_t = \mathbb{E}\|\bar m_t(\bar X_t) - m_t(\bar X_t)\|^2\), so
\eqref{eq:lem_reeves} reads
\(D_{\mathrm{KL}}(\bar P_0 \| P_0^\star) \le C_{\mathrm{diff}}\sum_t \Delta_t\),
which is the form used in Theorem~\ref{thm:kl}.

%=============================================================================
\section{Proofs for Section~\ref{sec:syndrome}}
\label{app:prop_mart}
\label{app:prop_syn}

%-----------------------------------------------------------------------------
\subsection{Proof of Proposition~\ref{prop:martingale}
            (Reverse-time martingale consistency)}

Recall \(M_t = m_t(X_t) = \mathbb{E}[X_0 \mid X_t]\) and \(X_t\) is obtained
from \(X_0\) through the Markov forward chain
\(X_0 \to X_1 \to \cdots \to X_T\) defined by the transition kernels
\(q(x_t\mid x_{t-1})\) (equivalently by
\(X_t = \alpha_t X_0 + \sigma_t\varepsilon_t\) with independent standard
Gaussian \(\varepsilon_t\)).

\begin{proof}[Proof of Proposition~\ref{prop:martingale}]
\textbf{Step 1 (Markov property).}
By construction of the forward chain, the joint law of \((X_0,X_{t-1},X_t)\)
factorizes as
\(q(x_0)\,q(x_{t-1}\mid x_0)\,q(x_t\mid x_{t-1})\). In particular, given
\(X_{t-1}\) the pair \((X_0, X_t)\) is conditionally independent:
\begin{equation}
X_0 \perp\!\!\!\perp X_t \;\big|\; X_{t-1}.
\label{eq:app_mart_ci}
\end{equation}

\textbf{Step 2 (Conditional mean reduction).}
By~\eqref{eq:app_mart_ci},
\(\mathbb{E}[X_0 \mid X_{t-1}, X_t] = \mathbb{E}[X_0 \mid X_{t-1}] = M_{t-1}\)
almost surely. Integrability of \(X_0\) (it is a square-integrable random
vector by assumption on \(P_0\)) together with boundedness of the Gaussian
transition kernels ensures that all conditional expectations exist and are
\(L^2\).

\textbf{Step 3 (Tower property).}
Applying the tower property with respect to the conditioning on
\((X_{t-1},X_t)\) and then on \(X_t\):
\begin{align*}
\mathbb{E}[M_{t-1} \mid X_t]
&= \mathbb{E}\!\left[\mathbb{E}[X_0\mid X_{t-1}] \,\big|\, X_t\right] \\
&= \mathbb{E}\!\left[\mathbb{E}[X_0 \mid X_{t-1}, X_t] \,\big|\, X_t\right]
  \qquad\text{(by Step 2)} \\
&= \mathbb{E}[X_0 \mid X_t]
  \qquad\qquad\text{(tower property)} \\
&= M_t.
\end{align*}
This holds \(P_t^\star\)-almost surely, proving the first form
of~Proposition~\ref{prop:martingale}. The second form is obtained by disintegration:
writing the above with the regular conditional distribution \(P(X_{t-1}\in
\cdot \mid X_t = x_t)\), which exists because the forward kernels are
Gaussian and hence Feller.

\textbf{Step 4 (Doob-martingale justification of terminology).}
Define the reverse-time filtration
\(\mathcal{G}_t = \sigma(X_t, X_{t+1}, \dots, X_T)\) for
\(t=0,1,\dots,T\), so \(\mathcal{G}_0 \supseteq \mathcal{G}_1 \supseteq \cdots
\supseteq \mathcal{G}_T\). By the forward Markov property, for any
\(t\in\{0,\dots,T-1\}\),
\(\sigma(X_0,\dots,X_t)\) is conditionally independent of
\(\sigma(X_{t+2},\dots,X_T)\) given \(X_{t+1}\); in particular,
\(\mathbb{E}[X_0 \mid \mathcal{G}_t] = \mathbb{E}[X_0 \mid X_t] = M_t\).
Hence \((M_t,\mathcal{G}_t)_{t=T}^{0}\) is precisely the Doob martingale of
\(X_0\) against the reverse filtration, and the backward-martingale property
\(\mathbb{E}[M_{t-1} \mid \mathcal{G}_t] = M_t\) specializes
to~Proposition~\ref{prop:martingale} when conditioning is reduced to \(X_t\). This
justifies calling \((M_t)\) a \emph{posterior-mean martingale}.
\end{proof}

%-----------------------------------------------------------------------------
\subsection{Proof of Proposition~\ref{prop:zero_syn}
            (Zero-syndrome property)}

\begin{proof}
Assume \(\tilde p_{t-1 \mid t}(\cdot\mid x_t) = p^\star_{t-1\mid t}(\cdot\mid
x_t)\) and \(\hat m_{t-1} = m_{t-1}\). For any measurable
\(f:\mathbb{R}^d\to\mathbb{R}^d\) with \(\mathbb{E}\|f(X_{t-1})\|<\infty\),
the operators defined by integration against these kernels coincide:
\((\tilde T_t f)(x_t) = (T_t^\star f)(x_t)\) for \(P_t^\star\)-a.e.\ \(x_t\).
Substituting \(f = \hat m_{t-1} = m_{t-1}\) and using the defect definition,
\begin{align*}
d_t(x_t)
&= \hat m_t(x_t) - (\tilde T_t \hat m_{t-1})(x_t)
\;=\; m_t(x_t) - (T_t^\star m_{t-1})(x_t).
\end{align*}
By Proposition~\ref{prop:martingale}, \(m_t = T_t^\star m_{t-1}\)
\(P_t^\star\)-a.e.\ (the pointwise form of~Proposition~\ref{prop:martingale}). Therefore
\(d_t(x_t) = 0\) \(P_t^\star\)-a.e. Since \(W_t\) is a deterministic matrix,
\(s_t(x_t) = W_t\,d_t(x_t) = 0\) on the same set.
\end{proof}

%-----------------------------------------------------------------------------
\subsection{Proof of Proposition~\ref{prop:proxy}
            (Unbiasedness of the one-sample proxy)}

\begin{proof}
Write
\(\hat s_t^{(1)} = W_t\bigl(\hat x_0(\tilde X_t,t) - \hat x_0(\tilde
X_{t-1}^{\mathrm{tent}},t-1)\bigr)\), where
\(\tilde X_{t-1}^{\mathrm{tent}} \sim \tilde p_{t-1\mid t}(\cdot \mid
\tilde X_t)\).

\textbf{Step 1 (Integrability for Fubini).}
By assumption on \(P_0\) and the Gaussian forward channel, \(\tilde X_t\)
and \(\tilde X_{t-1}^{\mathrm{tent}}\) are square-integrable, and under
(R1) the estimator \(\hat x_0(\cdot,t)\) is locally Lipschitz on the
relevant local ball, hence
\(\mathbb{E}\|\hat x_0(\tilde X_{t-1}^{\mathrm{tent}},t-1)\| < \infty\)
conditionally on \(\tilde X_t\). This integrability allows us to move
\(W_t\) and the conditional expectation through the defining integral.

\textbf{Step 2 (Conditional expectation of the lookahead term).}
Conditionally on \(\tilde X_t = x_t\), the distribution of \(\tilde
X_{t-1}^{\mathrm{tent}}\) is precisely \(\tilde p_{t-1\mid t}(\cdot\mid
x_t)\), so by definition of \(\tilde T_t\) applied to the vector-valued
function \(\hat m_{t-1} = \hat x_0(\cdot,t-1)\),
\begin{equation}
\mathbb{E}\!\left[\hat x_0(\tilde X_{t-1}^{\mathrm{tent}}, t-1)
  \,\big|\,\tilde X_t\right]
\;=\; (\tilde T_t \hat m_{t-1})(\tilde X_t).
\label{eq:app_proxy_step}
\end{equation}

\textbf{Step 3 (Assemble).}
Since \(W_t\) is deterministic and \(\hat x_0(\tilde X_t, t)\) is
\(\tilde X_t\)-measurable,
\begin{align*}
\mathbb{E}[\hat s_t^{(1)} \mid \tilde X_t]
&= W_t\!\left(\hat x_0(\tilde X_t,t) - \mathbb{E}[\hat x_0(\tilde
    X_{t-1}^{\mathrm{tent}},t-1) \mid \tilde X_t]\right) \\
&= W_t\!\left(\hat m_t(\tilde X_t) - (\tilde T_t\hat m_{t-1})(\tilde X_t)\right)
\;=\; s_t(\tilde X_t).
\end{align*}

For deterministic samplers, \(\tilde p_{t-1\mid t}(\cdot\mid x_t) =
\delta_{F_t(x_t;\omega_t)}\); the expectation in~\eqref{eq:app_proxy_step}
is a Dirac evaluation and the equality holds pointwise without averaging.
Hence \(\hat s_t^{(1)} = s_t(\tilde X_t)\) exactly.
\end{proof}

\paragraph{Variance bound for the \(K\)-sample proxy.}
Let \(\tilde X_{t-1}^{(1)},\dots,\tilde X_{t-1}^{(K)}\) be i.i.d.\ draws from
\(\tilde p_{t-1\mid t}(\cdot\mid\tilde X_t)\) and define
\begin{equation}
\hat s_t^{(K)}
\;=\;
W_t\!\left(\hat x_0(\tilde X_t,t) - \tfrac1K\sum_{k=1}^K \hat x_0(\tilde
X_{t-1}^{(k)},t-1)\right).
\label{eq:app_Ksample}
\end{equation}

\begin{lemma}[Variance of the \(K\)-sample proxy]\label{lem:Ksample}
Conditionally on \(\tilde X_t\),
\(\mathbb{E}[\hat s_t^{(K)} \mid \tilde X_t] = s_t(\tilde X_t)\) and
\begin{equation}
\mathbb{E}\!\left[\|\hat s_t^{(K)} - s_t(\tilde X_t)\|^2 \mid \tilde X_t\right]
\;=\;
\frac{1}{K}\,\mathrm{tr}\!\left(W_t\,\mathrm{Cov}_{\tilde p_{t-1\mid t}(\cdot\mid\tilde X_t)}[\hat x_0(\cdot,t-1)]\,W_t^\top\right).
\label{eq:Ksample_var}
\end{equation}
\end{lemma}

\begin{proof}
Unbiasedness is identical to Proposition~\ref{prop:proxy} since the sample
mean is an unbiased estimator of the conditional mean. For the variance,
let \(Y_k = \hat x_0(\tilde X_{t-1}^{(k)}, t-1)\) and
\(\mu = \mathbb{E}[Y_1 \mid \tilde X_t] = (\tilde T_t\hat m_{t-1})(\tilde
X_t)\). Then
\(\hat s_t^{(K)} - s_t(\tilde X_t) = -W_t\bigl(\tfrac1K\sum_k Y_k -
\mu\bigr)\). Since \(Y_1,\dots,Y_K\) are conditionally i.i.d.\ with
covariance \(\Sigma_Y = \mathrm{Cov}[Y_1 \mid \tilde X_t]\),
\(\tfrac1K\sum_k Y_k\) has conditional covariance \(\Sigma_Y/K\). Therefore
\begin{align*}
\mathbb{E}\!\left[\|\hat s_t^{(K)} - s_t(\tilde X_t)\|^2 \mid \tilde X_t\right]
&= \mathrm{tr}\!\left(W_t \cdot \Sigma_Y/K \cdot W_t^\top\right)
\;=\; \tfrac1K\,\mathrm{tr}(W_t \Sigma_Y W_t^\top).
\end{align*}
\end{proof}

Lemma~\ref{lem:Ksample} shows that the \(K\)-sample proxy reduces syndrome
estimator variance by a factor \(K\); for deterministic samplers \(\Sigma_Y=0\)
and all \(K\ge 1\) give the exact syndrome.

%=============================================================================
\section{Proofs for Section~\ref{sec:theory}}
\label{app:thm_proofs}

Throughout this section we work under (R1)--(R4) and under the local error
decomposition \(e_t = e_t^\parallel + e_t^\perp\) of the main text
(Section~\ref{sec:theory}). Operator norms \(\|\cdot\|\) on matrices are
spectral (largest singular value).

%-----------------------------------------------------------------------------
\subsection{Proof of Theorem~\ref{thm:observability}
            (First-order syndrome expansion)}
\label{app:thm_obs}

\begin{proof}
Let \(\bar x_t \in B(x_t^\star, r_t)\) and \(e_t = \bar x_t - x_t^\star\).
We decompose the martingale defect along the deployed process as
\begin{equation}
d_t(\bar x_t)
\;=\;
\underbrace{d_t(x_t^\star)}_{=:\,b_t^{(0)}}
\;+\;
\underbrace{\bigl(d_t(\bar x_t) - d_t(x_t^\star)\bigr)}_{=:\,\Delta_t^d}.
\label{eq:app_obs_split}
\end{equation}

\textbf{Step 1 (Bias term \(b_t^{(0)}\)).}
Using \(\hat m_t = m_t + (\hat m_t - m_t)\) and
\(\tilde T_t = T_t^\star + (\tilde T_t - T_t^\star)\),
\begin{align}
d_t(x_t^\star)
&= \hat m_t(x_t^\star) - (\tilde T_t \hat m_{t-1})(x_t^\star) \nonumber\\
&= \underbrace{\hat m_t(x_t^\star) - m_t(x_t^\star)}_{\text{(I)}}
 + \underbrace{m_t(x_t^\star) - (T_t^\star m_{t-1})(x_t^\star)}_{\text{(II)}=\,0} \nonumber\\
&\quad
 + \underbrace{(T_t^\star m_{t-1})(x_t^\star) - (T_t^\star \hat m_{t-1})(x_t^\star)}_{\text{(III)}}
 + \underbrace{(T_t^\star \hat m_{t-1})(x_t^\star) - (\tilde T_t \hat m_{t-1})(x_t^\star)}_{\text{(IV)}}.
\label{eq:app_obs_four_terms}
\end{align}
Term (II) vanishes by Proposition~\ref{prop:martingale}. By (R1) and the
definition of \(\zeta_t\), \(\|\text{(I)}\| \le \zeta_t\). Term (III) is
\(-T_t^\star(\hat m_{t-1} - m_{t-1})(x_t^\star)\); since \(T_t^\star\) is an
integration operator against a probability kernel, it is a contraction in
\(\sup\)-norm, so
\(\|\text{(III)}\| \le \|\hat m_{t-1}-m_{t-1}\|_\infty \le \zeta_{t-1}\).
Term (IV) is \((\tilde T_t - T_t^\star)\hat m_{t-1}\) evaluated at
\(x_t^\star\); under (R1) and the \(C^1\)-mismatch
\(\|\Psi_t - \Psi_t^\star\|_{C^1} \le \xi_t\), total-variation distance between
\(\tilde p_{t-1\mid t}\) and \(p^\star_{t-1\mid t}\) is
\(O(\xi_t)\) locally, so \(\|\text{(IV)}\| \le L_{m}\,\xi_t\) with
\(L_m = \|\hat m_{t-1}\|_{C^0(B)}\) finite by (R1). Collecting:
\begin{equation}
\|b_t^{(0)}\| \;\le\; (1+1)\zeta_t + L_m\,\xi_t
   \;\le\; c_{t,1}^{(0)}(\xi_t + \zeta_t),
\label{eq:app_obs_b0_bound}
\end{equation}
with \(c_{t,1}^{(0)} = \max(2, L_m)\). (For notational uniformity we
identify \(\zeta_{t-1}\) with the worst of \(\zeta_{t-1},\zeta_t\); under
(R1) both are uniformly bounded by a single \(\zeta_{\max}\), absorbed into
\(c_{t,1}^{(0)}\).)

\textbf{Step 2 (Linear term \(\Delta_t^d\) via Taylor's theorem).}
Under (R1), \(d_t\) is twice continuously differentiable on
\(B(x_t^\star,r_t)\) with bounded second derivatives. By Taylor's theorem,
\begin{equation}
d_t(\bar x_t) - d_t(x_t^\star)
\;=\; G_t e_t \;+\; R_t^{(2)}(e_t),
\label{eq:app_obs_taylor}
\end{equation}
where \(G_t = \nabla d_t(x_t^\star)\) and, with \(M_2\) a uniform bound on
\(\|\nabla^2 d_t\|\) over the ball,
\(\|R_t^{(2)}(e_t)\| \le \tfrac{M_2}{2}\|e_t\|^2\). Strictly speaking
\(\nabla d_t\) at \(x_t^\star\) equals
\(\nabla\hat m_t - \nabla(\tilde T_t \hat m_{t-1})\); writing
\(\nabla\hat m_t = \nabla m_t + (\nabla\hat m_t - \nabla m_t)\) and
similarly expanding the second term produces an extra
\((\xi_t+\zeta_t)\)-order contribution to \(G_t\) itself. We absorb this
into the first-order ``expansion around the exact process'':
\begin{equation}
d_t(\bar x_t) = G_t\, e_t \,+\, R_t^{(1)}(e_t),
\qquad
\|R_t^{(1)}(e_t)\| \le c'(\zeta_t + \xi_t)\|e_t\| + \tfrac{M_2}{2}\|e_t\|^2,
\label{eq:app_obs_Rone}
\end{equation}
where \(c'\) depends on \(\|\nabla^2 m_t\|\), \(\|\nabla \hat m_t\|\),
and the \(C^1\)-mismatch norms in (R1).

\textbf{Step 3 (Apply \(W_t\) and aggregate).}
Combining~\eqref{eq:app_obs_split},~\eqref{eq:app_obs_b0_bound},
and~\eqref{eq:app_obs_Rone}, and left-multiplying by \(W_t\),
\begin{equation}
s_t(\bar x_t) = W_t d_t(\bar x_t) = W_t G_t e_t + b_t + r_t,
\label{eq:app_obs_expansion_final}
\end{equation}
with \(b_t = W_t b_t^{(0)}\) and \(r_t = W_t R_t^{(1)}(e_t)\). Let
\(\|W_t\|_{\mathrm{op}} \le \sqrt{\overline w}\) (Assumption~\ref{asm:op}
in the main text; also recorded here). Then
\(\|b_t\| \le \sqrt{\overline w}\,c_{t,1}^{(0)}(\xi_t+\zeta_t) =:
c_{t,1}(\xi_t+\zeta_t)\) and
\(\|r_t\| \le \sqrt{\overline w}\bigl(c'(\xi_t+\zeta_t)\|e_t\|
+ \tfrac{M_2}{2}\|e_t\|^2\bigr)\). Using \(ab \le a^2/2 + b^2/2\) to expand
the cross term,
\(\|r_t\| \le c_{t,2}\bigl(\|e_t\|^2 + (\xi_t+\zeta_t)\|e_t\| +
(\xi_t+\zeta_t)^2\bigr)\) for a constant \(c_{t,2}\) depending on
\(\overline w, c', M_2\). This proves~\eqref{eq:syn_expansion} with the
claimed constants.

\textbf{Step 4 (Observability bound~\eqref{eq:obs_bound}).}
Decompose \(e_t = e_t^\parallel + e_t^\perp\) orthogonally; since
\(W_t G_t\) is linear,
\(W_t G_t e_t = W_t G_t e_t^\parallel + W_t G_t e_t^\perp\). By
Assumption~\ref{asm:obs}, \(\|W_t G_t e_t^\parallel\| \le \beta_t
\|e_t^\parallel\|\) and \(\|W_t G_t e_t^\perp\| \ge \kappa_t
\|e_t^\perp\|\). Rearranging the triangle inequality
applied to~\eqref{eq:app_obs_expansion_final},
\begin{align*}
\|s_t(\bar x_t)\|
&\;\ge\; \|W_t G_t e_t^\perp\| - \|W_t G_t e_t^\parallel\| - \|b_t\| - \|r_t\| \\
&\;\ge\; \kappa_t \|e_t^\perp\| - \beta_t \|e_t^\parallel\| - \|b_t\| - \|r_t\|,
\end{align*}
which gives
\(\|e_t^\perp\| \le \kappa_t^{-1}\bigl(\|s_t(\bar x_t)\| + \beta_t
\|e_t^\parallel\| + \|b_t\| + \|r_t\|\bigr)\), i.e.~\eqref{eq:obs_bound}.
The reverse triangle inequality is valid here because
\(\kappa_t > \beta_t\) by Assumption~\ref{asm:obs}, so the right side is
nonnegative once \(b_t, r_t\) are small enough.
\end{proof}

%-----------------------------------------------------------------------------
\subsection{Proof of Theorem~\ref{thm:stabilization}
            (Selective stabilization)}
\label{app:thm_stab}

We state the proof in the isotropic regime \(W_t=\Sigma_t=\Sigma_{t-1}=I_d\);
the general weighted case follows by the same argument after preconditioning
by \(W_t\) and the quadratic weights. Under Algorithm~\ref{alg:siec} with
\(L_t=1\), the corrected clean estimate is
\(\bar x_{0,t} = \hat x_{0,t} - \gamma_t \hat s_t\) and the corrected next
state is \(\bar x_{t-1} = \Psi_t(\bar x_t, \bar x_{0,t};\omega_t)\).

\begin{proof}
\textbf{Step 1 (Linearization of \(\Psi_t\) around the ideal trajectory).}
Under (R1) and the fact that \(\omega_t\) is frozen, the map
\((x,z)\mapsto\Psi_t(x,z;\omega_t)\) is \(C^2\) locally. Taylor-expanding
around \((x_t^\star, m_t(x_t^\star))\):
\begin{align}
\Psi_t(\bar x_t, \bar x_{0,t};\omega_t)
&= \Psi_t(x_t^\star, m_t(x_t^\star);\omega_t) \nonumber\\
&\quad+ \nabla_x \Psi_t^\star (x_t^\star,m_t(x_t^\star))\,e_t
 + \nabla_z \Psi_t^\star (x_t^\star,m_t(x_t^\star))\,(\bar x_{0,t} - m_t(x_t^\star))
 \nonumber\\
&\quad+ O(\|e_t\|^2 + \|\bar x_{0,t} - m_t(x_t^\star)\|^2 + \xi_t).
\label{eq:app_stab_Psi_taylor}
\end{align}
Here we have replaced \(\Psi_t\) by \(\Psi_t^\star\) in the linear terms
with a bookkeeping error of size \(\xi_t\), using the \(C^1\)-mismatch
bound in (R1). Set \(B_t = \nabla_x\Psi_t^\star\) and
\(C_t = \nabla_z\Psi_t^\star\) evaluated at \((x_t^\star, m_t(x_t^\star))\).
The zeroth-order term is
\(\Psi_t(x_t^\star,m_t(x_t^\star);\omega_t)\), which under (R1) equals
\(x_{t-1}^\star\) plus an \(O(\xi_t)\) deployment-mismatch contribution.

\textbf{Step 2 (Clean-estimate expansion).}
By Theorem~\ref{thm:observability} and the corrected clean estimate rule,
\begin{align*}
\bar x_{0,t} - m_t(x_t^\star)
&= \hat x_{0,t} - m_t(x_t^\star) - \gamma_t \hat s_t \\
&= \bigl(\hat m_t(\bar x_t) - m_t(x_t^\star)\bigr) - \gamma_t\bigl(G_t e_t + b_t + r_t\bigr).
\end{align*}
Expanding \(\hat m_t(\bar x_t) = m_t(x_t^\star) + J_t e_t + O(\zeta_t,
\|e_t\|^2)\) gives
\begin{equation}
\bar x_{0,t} - m_t(x_t^\star)
\;=\; J_t e_t - \gamma_t G_t e_t + \rho_t^{(z)},
\label{eq:app_stab_zbar}
\end{equation}
with
\(\|\rho_t^{(z)}\| \le c''(\zeta_t + \xi_t + \|e_t\|^2)\) for some
\(c''>0\).

\textbf{Step 3 (One-step corrected propagation).}
Substituting~\eqref{eq:app_stab_zbar} into~\eqref{eq:app_stab_Psi_taylor}
and subtracting the ideal update
\(x_{t-1}^\star = \Psi_t^\star(x_t^\star,m_t(x_t^\star);\omega_t)\) yields
\begin{equation}
e_{t-1}^{\mathrm{corr}}
\;=\; B_t e_t + C_t\bigl(J_t e_t - \gamma_t G_t e_t\bigr)
 + \rho_t^{(e)}
\;=\; L_t(\gamma_t)\,e_t + \rho_t^{(e)},
\label{eq:app_stab_onestep}
\end{equation}
where \(L_t(\gamma_t) = B_t + C_t J_t - \gamma_t C_t G_t\) as
in~\eqref{eq:L_operator}, and the residual satisfies
\(\|\rho_t^{(e)}\| \le c'''(\xi_t + \zeta_t + \|e_t\|^2)\) with
\(c'''\) depending on \(\|C_t\|\), \(c''\), and the second-derivative
bound from (R1). This is the key structural identity: the consensus update
injects the term \(-\gamma_t C_t G_t e_t\), exactly the term missing from
the uncorrected deployment.

\textbf{Step 4 (Projecting and applying Assumption~\ref{asm:decode}).}
Apply \(P_{t-1}^\perp\) to both sides of~\eqref{eq:app_stab_onestep} and
use \(I = P_t^\parallel + P_t^\perp\):
\begin{align*}
\|P_{t-1}^\perp e_{t-1}^{\mathrm{corr}}\|
&\le \|P_{t-1}^\perp L_t(\gamma_t) P_t^\perp e_t\|
   + \|P_{t-1}^\perp L_t(\gamma_t) P_t^\parallel e_t\|
   + \|P_{t-1}^\perp \rho_t^{(e)}\| \\
&\le \rho_t \|e_t^\perp\| + \chi_t \|e_t^\parallel\|
   + c_{t,3}(\xi_t + \zeta_t + \|e_t\|^2),
\end{align*}
using the norms in Assumption~\ref{asm:decode} and
\(\|P_{t-1}^\perp\| \le 1\), \(\|\rho_t^{(e)}\| \le c'''(\cdots)\), with
\(c_{t,3} = c'''\). This gives~\eqref{eq:stab_normal}.

Analogously for the tangent projection, using
\(\|P_{t-1}^\parallel L_t(\gamma_t) P_t^\parallel\| \le 1+\nu_t\),
\begin{align*}
\|P_{t-1}^\parallel e_{t-1}^{\mathrm{corr}}\|
&\le (1+\nu_t)\|e_t^\parallel\| + \chi_t \|e_t^\perp\|
   + c_{t,4}(\xi_t + \zeta_t + \|e_t\|^2).
\end{align*}
This is~\eqref{eq:stab_tangent}.

\textbf{Step 5 (Zero-correction degenerate case.)}
When \(\gamma_t = 0\), \(L_t(0) = B_t + C_t J_t\); without the
\(-\gamma_t C_t G_t\) term, the contraction bound
\(\|P_{t-1}^\perp L_t(0) P_t^\perp\|<1\) generally fails, since the ideal
one-step reverse map need not be a normal-contraction in this specific
decomposition. This recovers the accumulation regime
of~\citet{zhong2026iec}: without selective correction the normal component
grows, consistent with the exponential-to-linear gap they identify.
\end{proof}

%-----------------------------------------------------------------------------
\subsection{Proof of Theorem~\ref{thm:mmse} (Excess MMSE from syndrome energy)}
\label{app:thm_mmse}

\begin{proof}
\textbf{Step 1 (Core identity).}
Recall \(\bar d_t(x_t) = \bar m_t(x_t) - (\bar T_t\bar m_{t-1})(x_t)\),
\(\bar s_t = W_t\bar d_t\), \(m_t = T_t^\star m_{t-1}\) (by
Proposition~\ref{prop:martingale}), and
\(\Delta_t = \mathbb{E}\|\bar m_t(\bar X_t) - m_t(\bar X_t)\|^2\). Then
\begin{align}
\bar m_t - m_t
&= \bar m_t - T_t^\star m_{t-1} \nonumber\\
&= \bigl(\bar m_t - \bar T_t \bar m_{t-1}\bigr)
 + \bar T_t\bigl(\bar m_{t-1} - m_{t-1}\bigr)
 + \bigl(\bar T_t - T_t^\star\bigr) m_{t-1} \nonumber\\
&= \bar d_t + \bar T_t(\bar m_{t-1} - m_{t-1}) + (\bar T_t - T_t^\star)m_{t-1}.
\label{eq:app_mmse_id}
\end{align}

\textbf{Step 2 (Squared norm and three-term bound).}
Evaluate both sides of~\eqref{eq:app_mmse_id} at \(\bar X_t\), take squared
Euclidean norm, and apply
\(\|u+v+w\|^2 \le 3(\|u\|^2 + \|v\|^2 + \|w\|^2)\) (the standard bound from
Jensen applied to \(x\mapsto x^2\), equivalent to Cauchy--Schwarz with
weights \(1/3\)):
\begin{equation}
\|\bar m_t(\bar X_t) - m_t(\bar X_t)\|^2
\;\le\; 3\|\bar d_t(\bar X_t)\|^2
 + 3\|(\bar T_t)(\bar m_{t-1}-m_{t-1})(\bar X_t)\|^2
 + 3\|((\bar T_t - T_t^\star)m_{t-1})(\bar X_t)\|^2.
\label{eq:app_mmse_sq}
\end{equation}

\textbf{Step 3 (Operator stability on the middle term).}
Take expectation in~\eqref{eq:app_mmse_sq}. By Assumption~\ref{asm:op}
applied with \(f = \bar m_{t-1} - m_{t-1}\) (measurable, square-integrable
under (R1)--(R4)),
\(\mathbb{E}\|(\bar T_t)(\bar m_{t-1}-m_{t-1})(\bar X_t)\|^2
\le a\,\mathbb{E}\|(\bar m_{t-1}-m_{t-1})(\bar X_{t-1})\|^2 = a\,\Delta_{t-1}\).

\textbf{Step 4 (Weight conditioning on the defect term).}
\(\|\bar d_t\|^2 \le \underline w^{-1}\|W_t\bar d_t\|^2 = \underline
w^{-1}\|\bar s_t\|^2\) by Assumption~\ref{asm:op}(77) (every non-negative
quadratic form \(\|W_t u\|^2 = u^\top W_t^\top W_t u \ge \underline w \|u\|^2\),
so \(\|u\|^2 \le \underline w^{-1}\|W_t u\|^2\)). Taking expectations,
\(\mathbb{E}\|\bar d_t(\bar X_t)\|^2 \le \underline w^{-1}
\mathbb{E}\|\bar s_t(\bar X_t)\|^2\).

\textbf{Step 5 (Combine).}
Substituting Steps 3--4 into the expectation of~\eqref{eq:app_mmse_sq} and
using the definition
\(\eta_t^2 = \mathbb{E}\|((\bar T_t - T_t^\star)m_{t-1})(\bar X_t)\|^2\),
\begin{equation}
\Delta_t \;\le\; \frac{3}{\underline w}\,\mathbb{E}\|\bar s_t(\bar X_t)\|^2
  + 3a\,\Delta_{t-1} + 3\eta_t^2.
\label{eq:app_mmse_rec}
\end{equation}
This is~\eqref{eq:mmse_recursion}.

\textbf{Step 6 (Unroll the recursion).}
Iterating~\eqref{eq:app_mmse_rec} with \(\Delta_0 = 0\) yields
\begin{equation}
\Delta_t
\;\le\;
\sum_{s=1}^{t}(3a)^{t-s}\,\left(\tfrac{3}{\underline w}\mathbb{E}\|\bar s_s(\bar X_s)\|^2 + 3\eta_s^2\right)\bigg/ 3
\;=\;
\sum_{s=1}^{t}a^{t-s}\!\left(\tfrac{3}{\underline w}\mathbb{E}\|\bar s_s\|^2 + 3\eta_s^2\right).
\label{eq:app_mmse_unroll}
\end{equation}
(Strictly: iterating \(u_t\le A u_{t-1} + B_t\) with \(u_0=0\) and \(A\in[0,1)\)
gives \(u_t \le \sum_{s=1}^t A^{t-s}B_s\); here \(A=3a\) (if \(3a<1\)) or we
use the alternative form \(\Delta_t \le \sum_{s=1}^t a^{t-s} B_s'\) with
\(B_s' = 3\mathbb{E}\|\bar s_s\|^2/\underline w + 3\eta_s^2\), which is a
tighter geometric series under \(a\in[0,1)\).)

Summing~\eqref{eq:app_mmse_unroll} over \(t=1,\dots,T\) and switching the
order of summation,
\begin{align*}
\sum_{t=1}^T \Delta_t
&\;\le\;
\sum_{t=1}^T \sum_{s=1}^t a^{t-s}\!\left(\tfrac{3}{\underline w}\mathbb{E}\|\bar s_s\|^2 + 3\eta_s^2\right) \\
&= \sum_{s=1}^T\!\left(\tfrac{3}{\underline w}\mathbb{E}\|\bar s_s\|^2 + 3\eta_s^2\right)\sum_{t=s}^{T}a^{t-s}
\;\le\; \frac{1}{1-a}\sum_{s=1}^T\!\left(\tfrac{3}{\underline w}\mathbb{E}\|\bar s_s\|^2 + 3\eta_s^2\right),
\end{align*}
using \(\sum_{k=0}^{\infty}a^k = 1/(1-a)\). With
\(\mathcal{E}^{\mathrm{corr}}_{\mathrm{syn}} = \sum_t
\mathbb{E}\|\bar s_t(\bar X_t)\|^2\),
\begin{equation}
\sum_{t=1}^T \Delta_t
\;\le\;
\frac{3}{(1-a)\underline w}\,\mathcal{E}^{\mathrm{corr}}_{\mathrm{syn}}
 + \frac{3}{1-a}\sum_{t=1}^T \eta_t^2,
\label{eq:app_mmse_sum}
\end{equation}
which is~\eqref{eq:mmse_sum}.
\end{proof}

\paragraph{Remark (tightening constant 3 \(\to\) 2).}
The constant \(3\) in~\eqref{eq:app_mmse_sq} comes from the naive
\(\|u+v+w\|^2\le 3(\cdots)\). A Cauchy--Schwarz-type improvement with
weights \((\lambda_1,\lambda_2,\lambda_3)\) yields
\(\|u+v+w\|^2 \le (1+\sum_{i\ne 1}\lambda_i^{-1})\|u\|^2 + \cdots\); if one
is willing to handle \(\bar d_t\) with a larger weight and the two other
terms with smaller weights, the leading constant on
\(\mathcal{E}^{\mathrm{corr}}_{\mathrm{syn}}\) can be reduced to
\(\tfrac{2}{(1-a)\underline w}\) at the cost of inflating the \(\eta_t^2\)
coefficient; see Appendix~\ref{app:extra} for the bookkeeping.

%-----------------------------------------------------------------------------
\subsection{Proof of Theorem~\ref{thm:kl} (Distributional control)}
\label{app:thm_kl}

\begin{proof}
\textbf{Step 1 (Apply the comparison lemma).}
Under (R1)--(R4), Lemma~\ref{lem:reeves} gives
\begin{equation}
D_{\mathrm{KL}}(P_0^{\mathrm{corr}} \| P_0^\star)
\;\le\; C_{\mathrm{diff}} \sum_{t=1}^T
\mathbb{E}\|\bar m_t(\bar X_t) - m_t(\bar X_t)\|^2
\;=\; C_{\mathrm{diff}}\sum_{t=1}^T \Delta_t,
\label{eq:app_kl_step1}
\end{equation}
where \(C_{\mathrm{diff}}\) depends only on the schedule.

\textbf{Step 2 (Substitute Theorem~\ref{thm:mmse}).}
Substituting~\eqref{eq:app_mmse_sum} into~\eqref{eq:app_kl_step1},
\begin{equation}
D_{\mathrm{KL}}(P_0^{\mathrm{corr}} \| P_0^\star)
\;\le\; C_{\mathrm{diff}}\!
  \left[\frac{3}{(1-a)\underline w}\,\mathcal{E}^{\mathrm{corr}}_{\mathrm{syn}}
  + \frac{3}{1-a}\sum_{t=1}^T \eta_t^2\right]
= \frac{3C_{\mathrm{diff}}}{(1-a)\underline w}\,\mathcal{E}^{\mathrm{corr}}_{\mathrm{syn}}
 + \frac{3C_{\mathrm{diff}}}{1-a}\sum_{t=1}^T \eta_t^2,
\label{eq:app_kl_final}
\end{equation}
which is~\eqref{eq:kl_bound}.

\textbf{Step 3 (Pinsker).}
Pinsker's inequality states that for any two probability measures on the
same measurable space,
\(\|P - Q\|_{\mathrm{TV}} \le \sqrt{\tfrac12 D_{\mathrm{KL}}(P\|Q)}\).
Applied to \(P = P_0^{\mathrm{corr}}\) and \(Q = P_0^\star\),
\(\|P_0^{\mathrm{corr}} - P_0^\star\|_{\mathrm{TV}} \le \sqrt{\tfrac12
D_{\mathrm{KL}}(P_0^{\mathrm{corr}} \| P_0^\star)}\). Substituting
\eqref{eq:app_kl_final} bounds total-variation error in terms of cumulative
syndrome energy as well.
\end{proof}

%-----------------------------------------------------------------------------
\subsection{Proof of Corollary~\ref{cor:tradeoff}
            (Event-triggered compute--quality trade-off)}
\label{app:cor_tradeoff}

\begin{proof}
\textbf{Step 1 (Split accepted vs.\ triggered steps).}
Decompose
\begin{align*}
\mathcal{E}^{\mathrm{corr}}_{\mathrm{syn}}
&= \sum_{t=1}^T \mathbb{E}\|\bar s_t(\bar X_t)\|^2 \\
&= \sum_{t=1}^T \mathbb{E}\!\left[\|\bar s_t\|^2 \mathbf{1}\{r_t^{(0)}\le\tau_t\}\right]
 + \sum_{t=1}^T \mathbb{E}\!\left[\|\bar s_t\|^2 \mathbf{1}\{r_t^{(0)}>\tau_t\}\right].
\end{align*}

\textbf{Step 2 (Accepted steps).}
By Algorithm~\ref{alg:siec}, on accepted steps the corrected syndrome
\emph{equals} the tentative syndrome (no correction is applied), so
\(\|\bar s_t\|^2 = \|\hat s_t^{(0)}\|^2 = d\cdot r_t^{(0)}\). The trigger
rule enforces \(r_t^{(0)} \le \tau_t\), hence
\(\|\bar s_t\|^2 \le d\tau_t\) on accepted events. Therefore
\begin{equation}
\sum_{t=1}^T \mathbb{E}\!\left[\|\bar s_t\|^2 \mathbf{1}\{r_t^{(0)}\le\tau_t\}\right]
\le \sum_{t=1}^T d\tau_t\,\mathbb{P}(r_t^{(0)}\le\tau_t)
\le d\sum_{t=1}^T \tau_t.
\label{eq:app_trig_accepted}
\end{equation}

\textbf{Step 3 (Triggered steps).}
By the assumption of the corollary,
\(\mathbb{E}[\|\bar s_t\|^2 \mid r_t^{(0)}>\tau_t] \le q_t \mathbb{E}[\|\hat
s_t^{(0)}\|^2 \mid r_t^{(0)}>\tau_t]\). Multiplying both sides by
\(\mathbb{P}(r_t^{(0)}>\tau_t)\) (the conditioning probability):
\begin{equation}
\mathbb{E}\!\left[\|\bar s_t\|^2\mathbf{1}\{r_t^{(0)}>\tau_t\}\right]
\;\le\; q_t\,\mathbb{E}\!\left[\|\hat s_t^{(0)}\|^2\mathbf{1}\{r_t^{(0)}>\tau_t\}\right].
\label{eq:app_trig_triggered}
\end{equation}

\textbf{Step 4 (Combine).}
Adding~\eqref{eq:app_trig_accepted}
and~\eqref{eq:app_trig_triggered},
\begin{equation}
\mathcal{E}^{\mathrm{corr}}_{\mathrm{syn}}
\;\le\; d\sum_{t=1}^T\tau_t
+ \sum_{t=1}^T q_t\,\mathbb{E}\!\left[\|\hat s_t^{(0)}\|^2\mathbf{1}\{r_t^{(0)}>\tau_t\}\right],
\label{eq:app_trig_energy}
\end{equation}
which is~\eqref{eq:triggered_energy}.

\textbf{Step 5 (NFE bound).}
Each triggered step incurs at most \(L_t\) additional inner-loop iterations,
each adding one model evaluation. Therefore, in expectation,
\(\mathrm{NFE}_{\mathrm{extra}} = \sum_{t=1}^T N_t\) with
\(N_t \le L_t\,\mathbf{1}\{r_t^{(0)}>\tau_t\}\), and taking expectation
gives \(\mathbb{E}[\mathrm{NFE}_{\mathrm{extra}}]
\le \sum_{t=1}^T L_t\,\mathbb{P}(r_t^{(0)}>\tau_t)\).
\end{proof}

%=============================================================================
\section{Experimental Derivations and Additional Results}
\label{app:exp_details}

%-----------------------------------------------------------------------------
\subsection{Closed-form posterior mean for the low-rank Gaussian model}

Let \(X_0 = UZ\) with \(Z\sim\mathcal{N}(0,\Lambda)\), \(U\in\mathbb{R}^{d\times k}\),
\(U^\top U = I_k\), and \(\Lambda\in\mathbb{R}^{k\times k}\) diagonal and
positive. The forward channel gives \(X_t = \alpha_t UZ + \sigma_t\varepsilon\)
with \(\varepsilon\sim\mathcal{N}(0,I_d)\) independent of \(Z\).

\begin{lemma}[Posterior mean of the low-rank Gaussian]
\label{lem:lowrank_posterior}
The posterior \(Z \mid X_t = x_t\) is Gaussian with mean
\(\mu_t = \Lambda\alpha_t(\alpha_t^2\Lambda + \sigma_t^2 I_k)^{-1}U^\top x_t\)
and covariance \((\Lambda^{-1} + \alpha_t^2\sigma_t^{-2}I_k)^{-1}\). Hence
\begin{equation}
m_t(x_t) \;=\; \mathbb{E}[X_0\mid X_t = x_t] = U\mu_t
\;=\; U\Lambda\alpha_t\bigl(\alpha_t^2\Lambda + \sigma_t^2 I_k\bigr)^{-1} U^\top x_t,
\label{eq:app_lowrank_pmean}
\end{equation}
which is eq.~(89) of the main text.
\end{lemma}

\begin{proof}
The joint \((Z, X_t)\) is Gaussian:
\(X_t = \alpha_t UZ + \sigma_t\varepsilon\), so
\(\mathrm{Cov}(Z, X_t) = \alpha_t\Lambda U^\top\) and
\(\mathrm{Cov}(X_t) = \alpha_t^2 U\Lambda U^\top + \sigma_t^2 I_d\). Using
the Gaussian conditioning formula
\(\mathbb{E}[Z\mid X_t] = \mathrm{Cov}(Z,X_t)\mathrm{Cov}(X_t)^{-1}X_t\):
write \(\mathrm{Cov}(X_t) = \sigma_t^2 I_d + \alpha_t^2 U\Lambda U^\top\) and
apply the Woodbury identity:
\begin{equation*}
\mathrm{Cov}(X_t)^{-1} = \sigma_t^{-2}I_d -
  \sigma_t^{-2}U\bigl(\Lambda^{-1} + \alpha_t^2\sigma_t^{-2}I_k\bigr)^{-1}U^\top \alpha_t^2\sigma_t^{-2}.
\end{equation*}
Multiplying by \(\alpha_t\Lambda U^\top\) on the left gives
\(\mathbb{E}[Z\mid X_t] = \alpha_t\Lambda U^\top\mathrm{Cov}(X_t)^{-1}X_t\);
after simplification using \(U^\top U = I_k\),
\(\mathbb{E}[Z\mid X_t] = \Lambda\alpha_t(\alpha_t^2\Lambda +
\sigma_t^2 I_k)^{-1}U^\top X_t\). Multiplying by \(U\) gives
\(m_t(x_t)\) in~\eqref{eq:app_lowrank_pmean}.
\end{proof}

\subsection{Tangent and normal subspaces}
From~\eqref{eq:app_lowrank_pmean}, \(\nabla m_t(x_t) = U\,\Lambda\alpha_t
(\alpha_t^2\Lambda + \sigma_t^2 I_k)^{-1}U^\top\), which has rank \(k\) and
range \(\mathrm{span}(U)\). Thus
\(\mathcal{N}_t = \mathrm{range}(\nabla m_t) = \mathrm{span}(U)\) and
\(\mathcal{T}_t = \ker \nabla m_t = \mathrm{span}(U)^\perp\), as stated in
Section~\ref{sec:experiments}. The tangent subspace has dimension \(d-k\)
(diversity-preserving directions; lie in the null-space of the posterior
mean's sensitivity) and the normal subspace has dimension \(k\)
(structure-carrying directions).

\subsection{Verification of Assumption~\ref{asm:decode} for the
            low-rank Gaussian model}

\begin{lemma}[Selective decodability of the low-rank Gaussian model]
\label{lem:decode_verify}
For the low-rank Gaussian model with DDIM sampler and isotropic \SIEC
(\(W_t=I_d\), \(\Sigma_t=\Sigma_{t-1}=I_d\), \(\gamma_t\in(0,1)\)), there
exist positive constants \(\rho_t<1\), \(\nu_t=0\), \(\chi_t=0\) such that
Assumption~\ref{asm:decode} holds. In particular, \(L_t(\gamma_t)\) is
block-diagonal with respect to \(\mathcal{T}_t\oplus\mathcal{N}_t\).
\end{lemma}

\begin{proof}
For DDIM, \(\Psi_t(x,z) = \alpha_{t-1}z + (\sigma_{t-1}/\sigma_t)(x -
\alpha_t z)\), giving \(B_t = (\sigma_{t-1}/\sigma_t)I_d\) and
\(C_t = (\alpha_{t-1} - \alpha_t\sigma_{t-1}/\sigma_t)I_d =:
c_t^{\mathrm{DDIM}} I_d\). From Lemma~\ref{lem:lowrank_posterior},
\(J_t = U\Lambda\alpha_t(\alpha_t^2\Lambda + \sigma_t^2 I_k)^{-1} U^\top\)
lives in \(\mathrm{span}(U)\); similarly \(G_t\) (being a linear combination
of \(\nabla \hat m_t\) and \(\nabla(\tilde T_t\hat m_{t-1})\)) lives in
\(\mathrm{span}(U)\) when \(\hat m_t\) is the exact posterior mean plus an
isotropic perturbation as in the main text. Therefore
\(C_t J_t = c_t^{\mathrm{DDIM}} J_t\) and
\(C_t G_t = c_t^{\mathrm{DDIM}} G_t\) both vanish on \(\mathcal{T}_t\) and
restrict to endomorphisms of \(\mathcal{N}_t\).

Consequently, \(L_t(\gamma_t) = (\sigma_{t-1}/\sigma_t)I_d + c_t^{\mathrm{DDIM}}
(J_t - \gamma_t G_t)\) is block-diagonal in
\(\mathcal{T}_t\oplus\mathcal{N}_t\) with:
\begin{itemize}[leftmargin=1.2em]
  \item Tangent block: \((\sigma_{t-1}/\sigma_t) I_{\mathcal{T}_t}\).
  \item Normal block: \((\sigma_{t-1}/\sigma_t)I_{\mathcal{N}_t}
    + c_t^{\mathrm{DDIM}}(J_t|_{\mathcal{N}_t} - \gamma_t G_t|_{\mathcal{N}_t})\).
\end{itemize}
Thus \(\chi_t = 0\) (no crosstalk). The tangent block has norm exactly
\(\sigma_{t-1}/\sigma_t\), which for variance-preserving schedules with
\(\sigma_t^2 = 1-\alpha_t^2\) increasing towards \(T\) implies
\(\sigma_{t-1}/\sigma_t \le 1\); choosing \(\nu_t = 0\) suffices. For the
normal block, \(J_t|_{\mathcal{N}_t}\) is invertible and we choose
\(\gamma_t\) so that \(\|(\sigma_{t-1}/\sigma_t)I + c_t^{\mathrm{DDIM}}(J_t
- \gamma_t G_t)\|_{\mathcal{N}_t} =: \rho_t < 1\); Proposition~\ref{prop:proxy}
and an argument analogous to Theorem~\ref{thm:observability} show this is
achievable by taking \(\gamma_t = \lambda_t/(1+\lambda_t)\) with \(\lambda_t
= c\sigma_t^2/(\alpha_t^2+\sigma_t^2)\) and \(c\) in the range
\([0.3, 2]\), matching the saturation regime observed empirically in
Figure~\ref{fig:stabilization}.
\end{proof}

Lemma~\ref{lem:decode_verify} shows that for the exact testbed used in the
experiments, Assumption~\ref{asm:decode} holds with \(\chi_t = \nu_t = 0\)
and \(\rho_t\) explicitly computable from the schedule coefficients, the
eigenvalues of \(\Lambda\), and the correction gain.

\subsection{Exact KL divergence between two Gaussian DDIM chains}

For two deterministic DDIM chains driven by the same Gaussian seed and
same \(\omega_t\) but different clean-signal estimators
\(\hat m_t^{(1)}, \hat m_t^{(2)}\), the marginals at \(t=0\) are Gaussian
with means \(\mu^{(i)}\) and covariances \(\Sigma^{(i)}\) obtained by
propagating through the affine DDIM map. The KL divergence between two
Gaussians \(\mathcal{N}(\mu_1,\Sigma_1), \mathcal{N}(\mu_2,\Sigma_2)\) on
\(\mathbb{R}^d\) has the closed form
\begin{equation}
D_{\mathrm{KL}}\!\left(\mathcal{N}(\mu_1,\Sigma_1)\,\|\,\mathcal{N}(\mu_2,\Sigma_2)\right)
= \tfrac12\!\left[\log\tfrac{|\Sigma_2|}{|\Sigma_1|} - d
  + \mathrm{tr}(\Sigma_2^{-1}\Sigma_1)
  + (\mu_2-\mu_1)^\top\Sigma_2^{-1}(\mu_2-\mu_1)\right],
\label{eq:app_gauss_kl}
\end{equation}
evaluated exactly from sample means and covariances of the two chains.
This is the KL reported in Figure~\ref{fig:mmse_kl}(b).

\begin{figure}[htbp]
\centering
\includegraphics[width=0.7\linewidth]{fig3_mmse_kl.png}
\caption{\textbf{Syndrome energy as global quality indicator
(Theorems~\ref{thm:mmse}--\ref{thm:kl}).} (a) cumulative syndrome energy
\(\mathcal{E}^{\mathrm{corr}}_{\mathrm{syn}}\) vs.\ cumulative excess MMSE
\(\sum_t\Delta_t\); (b) same vs.\ exact KL divergence \(D_{\mathrm{KL}}\).
Each point corresponds to one \((\eta,c,\tau)\) configuration in the
sweep; the clear monotone relationship validates the accounting bound
\eqref{eq:mmse_sum}--\eqref{eq:kl_bound}.}
\label{fig:mmse_kl}
\end{figure}

\subsection{Figure 5: full trajectory comparison across deployment scenarios}

The figure omitted from the main text (Section~\ref{sec:experiments}
robustness paragraph) compares timestep-resolved error profiles across six
deployment scenarios: mild/moderate/severe isotropic perturbations, 8-bit
and 4-bit uniform quantization, and SNR-dependent perturbation. \SIEC
consistently reduces per-step error; the benefit scales with perturbation
severity; and the correction effect is visible throughout the reverse
process, not only at \(t=0\).

\begin{figure}[htbp]
\centering
\includegraphics[width=0.7\linewidth]{fig5_trajectories.png}
\caption{\textbf{Full trajectory error comparison across six deployment
scenarios.} Uncorrected (dashed), naive always-on refinement (dash-dot),
and event-triggered \SIEC (solid). \SIEC consistently reduces error
throughout the reverse process, with larger improvements under stronger
perturbation.}
\label{fig:trajectories}
\end{figure}

\subsection{Table~1: full compute--quality comparison}

Table~\ref{tab:pareto} reports the full sweep referenced in
Section~\ref{sec:experiments} (Corollary~\ref{cor:tradeoff} paragraph) at
\(\eta=0.3\); the resulting Pareto front is visualized in
Figure~\ref{fig:pareto}. Event-triggered \SIEC traces a smooth curve
parameterized by the trigger threshold \(\tau\): lowering \(\tau\)
allocates more NFE to correction and reduces path error, while raising
\(\tau\) conserves compute at the cost of residual error. At the matched
NFE budget of \(\sim 240\), \SIEC(\(\tau\)=80th) attains the best
compute-matched MSE (0.0913) and path error (0.280), beating both the
random-trigger and uniform-periodic baselines.

\begin{figure}[htbp]
\centering
\includegraphics[width=0.7\linewidth]{fig4_pareto.png}
\caption{\textbf{Pareto compute--quality trade-off
(Corollary~\ref{cor:tradeoff}).} (a) NFE vs.\ final MSE; (b) NFE vs.\ path
error. Event-triggered \SIEC (orange) is parameterized by the trigger
threshold \(\tau\); lowering \(\tau\) allocates more compute to correction
and reduces path error. \(\eta=0.3\), \(d=1024\), \(k=32\), \(T=200\).}
\label{fig:pareto}
\end{figure}

\begin{table}[htbp]
\centering
\caption{Compute--quality comparison at \(\eta=0.3\). 200 trajectories each.
Bold: best at matched NFE.}
\label{tab:pareto}
\small
\begin{tabular}{lccc}
\toprule
\textbf{Method} & \textbf{NFE} & \textbf{Final MSE} & \textbf{Path Error}\\
\midrule
Uncorrected & 200 & 0.0922 & 0.288 \\
\SIEC (\(\tau\!=\!95\)) & 223 & 0.0918 & 0.286 \\
\SIEC (\(\tau\!=\!90\)) & 243 & 0.0919 & 0.284 \\
Random trigger (20\%)   & 240 & 0.0919 & 0.281 \\
Uniform periodic (\(k\!=\!5\)) & 240 & 0.0918 & 0.281 \\
\SIEC (\(\tau\!=\!80\)) & 279 & \textbf{0.0913} & \textbf{0.280} \\
\SIEC (\(\tau\!=\!70\)) & 319 & 0.0917 & 0.277 \\
\SIEC (\(\tau\!=\!50\)) & 399 & 0.0920 & 0.270 \\
Naive refinement (always) & 400 & 0.0917 & 0.252 \\
\SIEC (always-on)       & 600 & 0.0915 & 0.253 \\
\bottomrule
\end{tabular}
\end{table}

%=============================================================================
\section{Additional Remarks}
\label{app:extra}

%-----------------------------------------------------------------------------
\subsection{\texorpdfstring{Tightening the constant $3$ in Theorem~\ref{thm:mmse}}{Tightening the constant 3 in Theorem 6.3}}

The constant \(3\) in~\eqref{eq:mmse_recursion} comes from the naive
bound \(\|u+v+w\|^2 \le 3(\|u\|^2+\|v\|^2+\|w\|^2)\). A more careful
Cauchy--Schwarz with weights \(\lambda_1,\lambda_2,\lambda_3>0\) summing to
\(1\) gives
\begin{equation*}
\|u+v+w\|^2 \;\le\; \lambda_1^{-1}\|u\|^2 + \lambda_2^{-1}\|v\|^2 + \lambda_3^{-1}\|w\|^2,
\end{equation*}
(equality iff \(\|u\|/\lambda_1 = \|v\|/\lambda_2 = \|w\|/\lambda_3\)). If
one expects the middle term \(\bar T_t(\bar m_{t-1}-m_{t-1})\) to dominate,
choose \(\lambda_2 = 1/2\) and \(\lambda_1 = \lambda_3 = 1/4\); this gives
\(\|u+v+w\|^2 \le 4\|u\|^2 + 2\|v\|^2 + 4\|w\|^2\). The MMSE recursion
becomes \(\Delta_t \le (4/\underline w)\mathbb{E}\|\bar s_t\|^2 + 2a\Delta_{t-1} + 4\eta_t^2\),
which summed yields
\(\sum_t \Delta_t \le (4/((1-2a)\underline w))\,\mathcal{E}^{\mathrm{corr}}_{\mathrm{syn}}
+ (4/(1-2a))\sum_t\eta_t^2\); the coefficient on the syndrome-energy term is
\(4\) instead of \(3\), but the constant \(a\) is effectively replaced by
\(2a\), tightening whenever \(a\) is very small. In practice, the choice
is an engineering one depending on which term is believed to be dominant.

%-----------------------------------------------------------------------------
\subsection{One-step limit and consistency models}

In the extreme one-step regime (consistency models
\citep{song2023consistency}), \(T = 1\) and the adjacent-step redundancy
disappears: there is no \((t,t-1)\) pair with \(t\ge 2\), so
\(\hat x_0(\tilde X_t,t)\) and \(\hat x_0(\tilde X_{t-1}^{\mathrm{tent}},
t-1)\) collapse to the same quantity at \(t=1\), making
\(\hat s_1^{(1)} \equiv 0\) by definition. Our framework degenerates: the
syndrome carries no information about deployment error. Extensions that
preserve the verifier-driven structure in this regime therefore require
additional sources of redundancy---for example, multiple independent inner
branches, multiple architectural pathways (skip connections, multi-scale
features), or multi-seed ensembles---and defining check equations across
these redundant sources rather than across timesteps.

%-----------------------------------------------------------------------------
\subsection{\texorpdfstring{Relation to $L^2$-projection decoders}{Relation to L2-projection decoders}}

The isotropic consensus update~\eqref{eq:iso_consensus} is precisely the
\(L^2\)-projection of \(\hat x_{0,t}\) onto the affine subspace
\(\{z : \|z - \hat x_{0,t}\|^2 + \lambda_t\|z - \hat x_{0,t-1}\|^2 =
\mathrm{const}\}\), weighted by the regularization strength \(\lambda_t\).
In classical decoding terminology, this is a soft-decision minimum-distance
decoder over two redundant noisy replicas of the same latent clean sample.
The non-isotropic version
\(\bar x_{0,t} = (\Sigma_t^{-1}+\lambda_t\Sigma_{t-1}^{-1})^{-1}(\Sigma_t^{-1}
\hat x_{0,t} + \lambda_t\Sigma_{t-1}^{-1}\hat x_{0,t-1})\)
corresponds to a noise-covariance-aware minimum-distance decoder in the
same sense. This connection makes the ``decoder'' terminology used in the
main text precise rather than metaphorical.

\end{document}